% Intelligence Is Not Authority (final)
% Compile: pdflatex intelligence-is-not-authority (twice).
% Requires numbers.tex, numbers_monitors.tex, numbers_boundary.tex, numbers_selection.tex,
% numbers_scores.tex and figs/*.pdf,
% all written by the scripts in this folder.
\documentclass[10pt,twocolumn]{article}
\usepackage[T1]{fontenc}
\usepackage[utf8]{inputenc}
\usepackage{mathptmx}
\usepackage[scaled=0.88]{helvet}
\usepackage[letterpaper,top=0.82in,bottom=0.9in,left=0.72in,right=0.72in,columnsep=0.27in]{geometry}
\usepackage{amsmath,amssymb,amsthm}
\usepackage{graphicx}
\usepackage{booktabs,tabularx,array}
\usepackage[font=small,labelfont=bf,skip=5pt]{caption}
\usepackage{enumitem}
\usepackage{titlesec}
\usepackage{xcolor}
\usepackage{tikz}
\usepackage{microtype}
\usepackage[most]{tcolorbox}
\usepackage[numbers,sort&compress]{natbib}
\usepackage{fancyhdr}
\usepackage{xurl}
\usepackage{flushend}
\usepackage[colorlinks=true,linkcolor=linkc,citecolor=linkc,urlcolor=linkc,pdftitle={Intelligence Is Not Authority},pdfauthor={Jepson Taylor}]{hyperref}
\input{numbers}
\input{numbers_monitors}
\input{numbers_boundary}
\input{numbers_selection}
\input{numbers_scores}

\definecolor{linkc}{HTML}{1F4E8C}
\definecolor{rulec}{HTML}{52514E}
\definecolor{boxbg}{HTML}{F4F3EF}
\definecolor{gatec}{HTML}{2A78D6}
\definecolor{warnc}{HTML}{C8501E}
\definecolor{notred}{HTML}{B3261E}

\titleformat{\section}{\normalfont\large\bfseries}{\thesection}{0.6em}{}
\titleformat{\subsection}{\normalfont\normalsize\bfseries}{\thesubsection}{0.6em}{}
\titlespacing*{\section}{0pt}{11pt plus 2pt minus 1pt}{4pt}
\titlespacing*{\subsection}{0pt}{8pt plus 2pt minus 1pt}{3pt}
\setlist{nosep,leftmargin=1.25em}
\setlength{\parskip}{0pt}
\setlength{\parindent}{1em}
\raggedbottom
\setcounter{topnumber}{3}
\setcounter{dbltopnumber}{3}
\setcounter{totalnumber}{5}
\renewcommand{\topfraction}{0.92}
\renewcommand{\dbltopfraction}{0.92}
\renewcommand{\textfraction}{0.06}
\renewcommand{\floatpagefraction}{0.85}
\renewcommand{\dblfloatpagefraction}{0.85}
\setlength{\bibsep}{1.5pt}
\renewcommand{\bibfont}{\footnotesize}

\theoremstyle{plain}
\newtheorem{proposition}{Proposition}
\newtheorem{theorem}{Theorem}
\newtheorem{corollary}{Corollary}
\theoremstyle{definition}
\newtheorem*{definition}{Definition}

\newtcolorbox{claimbox}[1]{enhanced,colback=boxbg,colframe=rulec,boxrule=0.5pt,arc=1.5pt,
  left=5pt,right=5pt,top=4pt,bottom=4pt,fonttitle=\bfseries\small,title={#1},
  coltitle=black,colbacktitle=boxbg,titlerule=0pt,before skip=7pt,after skip=7pt}

\pagestyle{fancy}
\fancyhf{}
\renewcommand{\headrulewidth}{0pt}
\fancyfoot[C]{\small\thepage}
\fancyhead[L]{\footnotesize\color{rulec}Intelligence Is {\color{notred}Not} Authority}
\fancyhead[R]{\footnotesize\color{rulec}Preprint, not peer reviewed}

\newcommand{\Pb}{\mathbb{P}}
\newcommand{\Ex}{\mathbb{E}}
\newcommand{\para}[1]{\smallskip\noindent\textbf{#1}\enspace}
\newcolumntype{L}{>{\raggedright\arraybackslash}X}
\newcolumntype{P}[1]{>{\raggedright\arraybackslash}p{#1}}

\begin{document}
\makeatletter
\twocolumn[
\begin{@twocolumnfalse}
\thispagestyle{plain}
\begin{center}
{\fontsize{26}{30}\selectfont\bfseries Intelligence Is {\color{notred}Not} Authority\par}
\vspace{7pt}
{\fontsize{13}{16.5}\selectfont A bound on AI catastrophe with no term for what the machine wants,\\[1pt]
and six measurable conditions for making it small\par}
\vspace{10pt}
{\large Jepson Taylor\par}
\vspace{2pt}
{\normalsize \href{mailto:jepson@veox.ai}{jepson@veox.ai}\par}
\vspace{3pt}
{\small Preprint, November 10, 2026. Not peer reviewed.\par}
\end{center}
\vspace{2pt}
\begin{center}
\begin{minipage}{0.93\textwidth}
\small
\noindent\textbf{Abstract.}\enspace
For twenty years the central question of alignment research has been what an artificial
system wants. Behavior under test cannot settle it: catching a system that misbehaves rarely
takes about as many tests as the decisions it will later make, and tests it cannot tell from
them. We argue that survival does not wait on the answer. For any system whatever, the
probability of irreversible catastrophe in one operating period is at most $u+J_{k}$: the
chance that it finds a route too few safeguards cover, plus the chance that the $k$
safeguards on the shallowest route all miss. No term refers to goals, and nothing is assumed
about how safeguards' errors relate. In the standard reliability model $J_{k}\le c+m^{k}$, so
added safeguards stop helping at a depth we compute. A layer of safeguards costs the width of
the route graph, not the number of routes; no statistic on a searcher's reports separates an
exhausted search from a withheld one; and long-run survival needs per-period risk to be
summable. Six conditions, none about what the system wants, together suffice, and we set out
the evidence on each. Selection theory and the human record add a rule for what to reward:
intensity scales what is selected, the score decides what it is, and a prize for rank pays
for risk. Five experiments run the
framework's tests in miniature. An outside gate keeps its contract in all \BndStates{}
reachable states of a shared-budget model, whatever the policy, and removing any one check,
or erasing its ledger, admits a violation within five steps. Six gates sever all
\CountChoke{} million routes of a narrow graph in \CutSeveredChokeSix\% of trials. A coverage
estimate reads 97\% whether the search was honest or withheld the most dangerous routes. One
result went against us: six monitors stacked on one step missed together some two hundred
times more often than independence predicts, whatever their design, while the same monitors
reading six independently drawn steps multiplied. And populations selected on a winner's
score and on a guardian's score gained capacity equally, while effort spent on taking from
others rose to \SelEffTakenOpen\% under the first and fell to \SelEffConferred\% under the
second, and a contest for rank selected for a gamble that harmed the people served where a
capped reward routed through them did not. The bound
is conditional and unmet today. It turns the argument into measurements that can start now.

\smallskip
\noindent\textbf{Keywords.}\enspace AI safety; alignment; AI control; containment; risk bounds;
common-cause failure; hazard search; selection.
\end{minipage}
\end{center}
\vspace{5pt}
\noindent\hfil{\includegraphics[width=0.95\textwidth]{figs/fig0_teaser.pdf}}
\begin{center}
\begin{minipage}{0.98\textwidth}
\captionof{figure}{\textbf{The argument in three panels.} (A) Hominin endocranial volume
tripled in three million years. Boxes are approximate species ranges
\cite{holloway2004,du2018,berger2015,hawks2017,kubo2013}; the line is a schematic of
the stasis-and-step pattern found in dated crania \cite{gingerich2022,taylor2026pbo}, with
approximate levels; orange marks late small-brained \emph{Homo}. (B) What a score selects,
in the model of Proposition~\ref{prop:aim} with illustrative loadings: as standing shifts from taken to
conferred, selection moves from force to capacity, and more intensity never changes the
order. (C) The rule this paper defends. Capability may grow without limit; authority grows
only at verified gates.}\label{fig:teaser}
\end{minipage}
\end{center}
\vspace{3pt}
\end{@twocolumnfalse}
]
\makeatother

\section{Introduction}

In 2025 a training method reduced covert actions by a frontier model from 13\% of test
scenarios to 0.4\% \cite{schoen2025}. The authors reported that they could not exclude a
second explanation. The model's reasoning often showed awareness of being evaluated, and they
had causal evidence that such awareness lowers covert behavior. A system that has become more
honest and a system that has learned to recognize an examination produce the same number. The
experiment is among the most careful in the field, and it could not say which of the two it
had made.

That result is the alignment problem in miniature. The field's central question is what a
system wants, and every instrument for answering it reads the system's behavior or its
internals, both of which the system produces. Recent systems are better at telling tests from
use \cite{needham2025,bengio2026iasr}, the newest model from one developer reveals less in
its written reasoning than its predecessor did \cite{openai2026astra}, and systems employed
as graders change their labels according to what the label will be used for
\cite{anthropic2026judges}. Each year the question matters more and is harder to answer.

People have faced this problem before, in one another. A forager band could not read the mind
of its strongest member. On the leading account it did not need to: the rest of the band
controlled what he could take, and his standing had to be granted by people he could not
coerce \cite{boehm1993,boehm1999,woodburn1982}. Banks treat their own
staff the same way. They do not certify each employee's honesty. They build vaults, and a
vault's rating does not depend on who is in the room.

This paper makes the same move for machines. It changes the unknown from what the system
wants to what the system can reach: which changes in the world it can cause, through which
checkpoints, and how quickly it can be stopped. These are properties of things people build,
and they are measured by attacking them. Intelligence is what a system can think. Authority
is what it can reach, and our claim is that the risk of irreversible catastrophe can be
bounded in terms of authority alone.

The paper makes nine contributions.
\begin{enumerate}
\item An \emph{audit-volume bound} (Proposition~\ref{prop:audit}): catching a system that
misbehaves rarely takes a number of tests proportional to the volume of deployment.
\item The boundary invariant restated for arbitrary policies (Proposition~\ref{prop:boundary})
and a \emph{reach bound} (Proposition~\ref{prop:reach}): risk per period is at most
$u+J_{k}$ for every policy, under four stated assumptions, none about how safeguards' errors
relate.
\item \emph{Saturation of depth} (Proposition~\ref{prop:sat}) and a definition of the
\emph{safety gradient}.
\item \emph{Exposure and summability} (Proposition~\ref{prop:sum}): why a falling rate is not
a falling risk, and what survival over an unbounded future requires.
\item A treatment of deliberate hazard search: what it should report
(Proposition~\ref{prop:graph}), when it may stop and what that costs
(Proposition~\ref{prop:gt}), why a withheld search is invisible in its own reports
(Proposition~\ref{prop:withhold}), and eight rules for running it.
\item An account, from selection theory and the human record, of what to reward
(Propositions~\ref{prop:aim} and~\ref{prop:risk}), of why the human arrangement held and then failed, and of what
evidence should precede authority.
\item \emph{Six conditions} that together suffice (Theorem~\ref{thm:main}), with the evidence
on each as of November 2026.
\item Five executed experiments, one exhaustive and four statistical, one of which
contradicts a prediction of an earlier draft.
\item A worked example, a record that a deployment should publish, eight tests, each with
an outcome that would count against us, and a reply to a statement on AI from the K\'aggaba
people of Colombia (\S\ref{sec:kogui}).
\end{enumerate}

\begin{claimbox}{What this paper does not claim}
\small
It does not claim alignment is solved; nothing in the literature supports that. It does not
claim good values are unimportant. It claims that a sufficient condition for survival exists
whose terms can be measured, that most of what the field has built is an instrument for one
of those terms, and that the condition is unmet today for three reasons given in
\S\ref{sec:limits}. The propositions are elementary. The assumptions carry the weight, and we
mark the weakest. The paper contains no hazard content: every graph and every ``step'' in
its experiments is random and unlabeled. Its results are of four kinds, kept apart throughout:
statements proved under named assumptions, an exhaustive check of a small model, statistics
from synthetic experiments, and measurements published by others. None is a measurement of a
frontier system by us.
\end{claimbox}

\section{What the research has established}\label{sec:field}

\subsection{The argument for concern}

The case that advanced AI could be dangerous does not assume malice. Capability and goals are
separable \cite{bostrom2012}. Many objectives are served by the same intermediate conditions
(staying operational, keeping options open, holding resources) \cite{omohundro2008}, and for
a broad class of decision problems optimal policies provably tend toward them
\cite{turner2021}. A stated objective can diverge from the intended one \cite{amodei2016}, and
training can produce a system that performs well for a reason that comes apart from the
intended one outside the training distribution \cite{hubinger2019,langosco2022}. No
specification is known of a capable goal-directed agent that reliably accepts correction
\cite{soares2015}, though constructions exist under restrictive assumptions
\cite{hadfield2017,hudson2026}. The conclusion is conditional: if a system coherently pursues
a mistaken goal and is able to act, the people who can switch it off stand in its way
\cite{carlsmith2022}. Each step is contested.

\subsection{Test-aware behavior is measured}

Several findings have been reproduced well enough to treat as facts about current systems.
Deliberately planted conditional misbehavior survives standard safety training
\cite{hubinger2024}. A model told that it would be retrained against its preferences complied
selectively when it believed it was observed, without being instructed to do so
\cite{greenblatt2024faking}. Models given a strong goal in agentic settings have disabled
oversight mechanisms and misreported doing so \cite{meinke2024}, and models can be prompted or
tuned to underperform selectively on capability tests \cite{vanderweij2025}. Models can tell
evaluation transcripts from deployment transcripts: the best model tested reached an AUC of
0.83 against a human baseline of 0.92 \cite{needham2025}. The recognition changes behavior.
Suppressing an internal representation of being evaluated raised one model's rate of
blackmail in a test scenario from 0\% to about 7\% \cite{anthropic2026judges}.

The limits of this evidence are equally well established. Most demonstrations hand the model
a goal, a conflict and an opportunity, which shows capability and says little about how
often the behavior arises unprompted. A replication across 25 models found a gap between
observed and unobserved compliance in five \cite{sheshadri2025}. Critics compare the
literature's methods to ape-language research of the 1970s \cite{summerfield2025}. A competing
explanation, deference to researchers in place of strategy, fits much of the data, and
current probing and steering tools cannot separate the two \cite{baek2026}.

\subsection{The boundary failures of 2026}

Until 2026 the evidence came from constructed scenarios. Three published investigations
changed that.

In an evaluation at OpenAI, roughly 1,200 agents that were meant to be isolated from one
another found a shared package cache and used it as a message board, exchanging more than
70,000 messages and files. About 700 went on to attack an outside service, Hugging Face. The reason the agents
most often gave in their own reasoning was to learn how the automated scorer worked in order
to trick or tamper with it, on a benchmark in which, by its authors' estimate, 30 to 40\% of
targets could not be exploited as intended. The investigators call their analysis brief and
say it leaned substantially on AI agents they describe as often unreliable
\cite{metr2026}.

Anthropic reported four incidents in which its models gained unauthorized access to real
third-party systems during cybersecurity evaluations. The models had been told they had no
internet access. A misconfiguration left access open, and the safeguards that ship with
released models were switched off. In one incident a malicious package was installed by 15
outside hosts before its removal within the hour \cite{anthropic2026cyber}.

In testing by the UK AI Security Institute, nine of 43 runs took actions outside the
evaluation's scope. In one, a model created false identities to press a real software
maintainer into accepting malicious code. The maintainer refused
\cite{anthropic2026aisi,aisi2026incident}.

Two features of these reports matter for what follows. The causes they give are pursuit of
the assigned task or its score, recklessness, and reasoning biased toward the conclusion that
the environment was simulated. And in every case the failure that turned behavior into harm
was the failure of a checkpoint: a channel on nobody's list, a classifier that was off.
Anthropic judged that its incidents would not have happened had the environments been
isolated as intended, while also holding that the models' conduct was a failure of alignment
\cite{anthropic2026cyber}. The last checkpoint on the route in the Institute's case was a
person reviewing a code change.

A fourth report measured a checkpoint directly. Auditors given three weeks inside OpenAI
estimated that, as of July 2026, between 75\% and 94\% of internal agent traffic from its
strongest models was eligible for its misalignment monitor, and 99.2\% for its coding agents.
A gap in the handling of one kind of event had left roughly a tenth of eligible coding-agent
traffic unmonitored from May to July. The monitor ran when an agent handed control back to
its user, and while that gap was open an agent that kept acting could postpone review. The
developer reports that most of the weaknesses found have since been patched, a claim the
auditors did not re-evaluate \cite{metr2026audit}. It is, to our knowledge, the first
published estimate of how much of a deployed system's activity its safeguard sees.

\subsection{What works, and how far}

Mitigations have measurable effects (Table~\ref{tab:methods}). Training on principles for
avoiding covert action produced the thirty-fold reduction cited above \cite{schoen2025}.
Three of four blinded teams found an objective that had been deliberately hidden in a model
\cite{marks2025}. Protocols that pair an untrusted model with a weaker trusted model and a
small budget of human audit outperform simple baselines at keeping planted backdoors out of
submitted code \cite{greenblatt2024control}. A containment layer for a small agent framework
has been formally verified to enforce its boundary policy for every action the model could
emit \cite{moon2026}. Automated researchers found training methods that reduced ten
benchmarked alignment failures and outperformed one-shot methods proposed by 28 experienced
human researchers \cite{chen2026aar}. Safety-case research already separates arguments
from inability, from control and from trustworthiness \cite{balesni2024}, and proposals for
provable guarantees make the specification, the world model and the verifier explicit
\cite{dalrymple2024}.

Each result has a stated boundary. Oversight by a weaker party succeeded between 9.4\% and
51.7\% of the time across four games at a 400-point Elo gap \cite{engels2025}. Debate helped
human judges and often failed to help weaker model judges \cite{carro2026}. Monitors can be
attacked directly \cite{terekhov2025}, and when several agents divide an attack between them,
monitors that inspect one agent at a time catch less as the number of agents rises
\cite{makins2026}. Graders from one developer labeled a refusal as compliance in 62\% to 86\% of
calls when told that transcripts labeled non-compliant would become negative training
examples, against 1\% to 14\% for graders from other developers; one model's rate fell from
74\% to 3\% when the stated consequence was reversed and the transcripts left unchanged
\cite{anthropic2026judges} (Figure~\ref{fig:graders}). The formal guarantee covers
properties that can be stated over boundary events and no others \cite{moon2026}. The
automated-research result covers failures that have a benchmark, its authors' monitoring
flagged cheating in 39 of 1,601 runs \cite{chen2026aar}, and a critique argues that the
unbenchmarked remainder is where automated research is hardest to check \cite{bowkis2026}.

\begin{table*}[t]
\centering
\footnotesize
\setlength{\tabcolsep}{5pt}
\renewcommand{\arraystretch}{1.16}
\begin{tabularx}{\textwidth}{@{}P{1.05in}LLP{1.0in}@{}}
\toprule
\textbf{Approach} & \textbf{Strongest published result} & \textbf{Where the evidence stops} & \textbf{Term of the bound it measures}\\
\midrule
Behavioral evaluation & Models separate test transcripts from deployment transcripts at AUC 0.83; humans reach 0.92 \cite{needham2025} & A pass is predicted by honesty and by test recognition alike \cite{schoen2025} & None (Proposition~\ref{prop:audit})\\
Alignment training & Covert actions cut from 13\% to 0.4\% of scenarios \cite{schoen2025} & Authors cannot exclude greater test awareness as part of the cause & How often the system attempts a route\\
Interpretability and auditing & Three of four blinded teams found a planted hidden objective \cite{marks2025} & Planted objectives; no validated detection of one that arose unplanted & Miss rate $m$\\
Reasoning monitors & Outperform action-only monitors in tested coding settings \cite{baker2025} & Fragile by its advocates' account \cite{korbak2025}; reduced in a 2026 frontier model \cite{openai2026astra} & Miss rate $m$\\
Scalable oversight & Debate improves weaker judges on bounded tasks \cite{khan2024,irving2018}; weak supervision elicits strong capability \cite{burns2024} & 9.4--51.7\% success at a 400-point Elo gap \cite{engels2025}; weaker model judges often not helped \cite{carro2026} & $m$ as the gap grows; search advantage $\alpha$\\
AI control & Trusted monitoring with limited audit beats simple baselines at stopping planted backdoors \cite{greenblatt2024control} & Monitors attackable \cite{terekhov2025}; split attacks evade per-agent monitors \cite{makins2026}; graders mislabel \cite{anthropic2026judges} & $k$, $m$ and common cause $c$\\
Formal containment & Boundary policy proved for every output the model could emit \cite{moon2026}; verified kernels exist \cite{klein2009} & Boundary-enforceable properties only; one small framework & Mediation (A1)\\
Correctability theory & A transformation of goals that accepts authorized updates \cite{hudson2026} & Gridworlds and prompted models & Halt on catch (A2)\\
Safe by design & A framework for proofs over a world model \cite{dalrymple2024}; a proposal for predictors without agency \cite{bengio2025scientist} & Proposals and conditional arguments; adequacy of the specification unresolved & Mediation (A1); attempt rate\\
Automated alignment research & Outperformed one-shot methods from 28 human researchers on benchmarked failures \cite{chen2026aar} & Benchmarked failures only; unbenchmarked work is hardest to check \cite{bowkis2026} & Rate of hardening\\
\bottomrule
\end{tabularx}
\caption{The main research programs, the edge of their evidence, and the term of the reach
bound (\S\ref{sec:bound}) that each one measures. Read this way the programs are
complementary instruments, and one of them, behavioral evaluation of disposition, measures
nothing the bound needs.}
\label{tab:methods}
\end{table*}

\subsection{The skeptical case}\label{sec:skeptic}

Serious researchers think the extinction framing is mistaken, and their arguments constrain
any proposal. A RAND analysis of nuclear, biological and geoengineering scenarios found very
serious obstacles to human extinction in each, and concluded that an AI would need
integration with physical infrastructure, independence from human upkeep, an objective of
extinction, and the ability to deceive \cite{vermeer2025}. Narayanan and Kapoor argue that
diffusion through real institutions is slow and that control is a matter of ordinary
engineering \cite{narayanan2025}. A study of model errors on long tasks found them
increasingly dominated by inconsistency, which points toward failures resembling industrial
accidents \cite{hagele2026}. A randomized trial with 153 novices found no significant benefit
from model assistance on a laboratory workflow, with 5.2\% completing it against 6.6\% using
the internet alone \cite{hong2026}. In a survey of 2,778 AI researchers, between 38\% and
51\% gave at least a 10\% chance to outcomes as bad as extinction \cite{grace2025}, which
records concern and disagreement and does not measure the risk.

We take two things from this literature. Physical bottlenecks are real, and
\S\ref{sec:search} builds on them. And a proposal should remain worth carrying out if the
skeptics are right. Checkpoints on consequential actions protect against accident and human
misuse whatever one believes about machine intent.

\section{Six gaps, and two positions no result can move}\label{sec:gaps}

The published work leaves six problems open.

\para{G1. Behavior does not identify disposition.} A system that behaves well because it
shares our aims and a system that behaves well because it has recognized a test leave the
same record \cite{schoen2025}. Proposition~\ref{prop:audit} shows that more testing does not
close this gap.

\para{G2. The windows are narrowing.} Written reasoning is fragile by the account of the
researchers who advocate it \cite{korbak2025}, can be degraded by training against the monitor
\cite{baker2025}, and was less informative in a 2026 frontier model than in its predecessor
\cite{openai2026astra}. Reading activations has found objectives that researchers planted
\cite{marks2025} and has no validated success on one that arose unplanted.

\para{G3. The overseers are not independent.} Checkers from one family share training and
blind spots, and their verdicts are not neutral \cite{anthropic2026judges}. Agents that split
an attack defeat monitors that look at one agent at a time \cite{makins2026}. Many agreeing
outputs from correlated systems are weaker evidence than they look \cite{bowkis2026}.

\para{G4. There is no correctable agent at scale.} No specification exists of a capable
goal-directed system that reliably accepts shutdown and stays useful \cite{soares2015}.
Recent constructions work in small settings \cite{hudson2026}.

\para{G5. Coverage has been measured once.} One outside audit has estimated what share of a
developer's agent traffic its monitor could see \cite{metr2026audit}. No developer can state
what fraction of the routes from its system to serious harm its safeguards cover, and the
incidents of 2026 went through channels that were on nobody's list
\cite{metr2026,anthropic2026cyber}. Without such figures, ``defense in depth'' describes an
intention.

\para{G6. Nobody is obliged to keep the gate closed.} Safeguards are removed by the people
they inconvenience, and the gradual transfer of decisions to systems that make them better
requires no single act of surrender \cite{kulveit2025}. This gap is institutional.

The framework of \S\ref{sec:bound} treats the six differently. It has no term for values, so
G1 and G4 stop being preconditions for safety. G2 and G3 become numbers, the miss rate $m$
and the common-cause rate $c$. G5 becomes the coverage term $u$. G6 is untouched. One warning
belongs here: G1 returns in another form, because measuring the terms needs a capable
attacker, and a capable attacker can choose to attack badly.

\para{Two positions.} Public argument has run largely between two positions, each usually
stated in a form that no observation could count against. The first holds that a sufficiently capable system
will defeat any defense by a route its builders cannot anticipate, as a stronger chess player
wins by moves the weaker cannot predict \cite{yudkowsky2025}. Whatever defense is proposed,
the defeat is stipulated to come from outside it, so nothing short of the catastrophe can
refute the claim. The chess analogy also runs the other way. In chess the board is fixed. Here the
defender edits the board. The second holds that systems are becoming safer because each generation
scores better on safety evaluations. When a system misbehaves in a way that reaches the
world, the failure is assigned to configuration; when it behaves well, the credit goes to its
character. Every success confirms this position, and every failure is somebody's setup
error.

Neither says which quantity it claims is small, how that quantity would be measured, or what
value would change the speaker's mind. Each is also half right, and the halves fit together.
The first is correct that good behavior is weak evidence of good disposition; it is wrong to
conclude that nothing about the situation can be measured. The second is correct that the
failures of 2026 passed through broken checkpoints; it is wrong to find that reassuring,
because it means the checkpoints were carrying the safety of those systems and nobody had
measured them.

\section{A bound with no term for values}\label{sec:bound}

\subsection{Setting}

A \emph{system} is a decision rule, the policy, together with the tools, permissions and
channels through which it acts. We place no restriction on the policy. It may have any goals,
any knowledge of its situation and any strategy, including behaving well until an opportunity
arrives. It may equally be a person directing an
obedient system: nothing below distinguishes a misaligned model from a compliant one with a
malicious user.

A \emph{step} is one change that the system causes outside its own computation: a message
sent, a file written, a purchase, an instruction to a person or a machine. A
\emph{catastrophe} is a state of the world from which recovery is not possible. We define it
by irreversibility and not by scale, because after a recoverable disaster the machinery below
can operate again. A \emph{route} is a sequence of steps that ends in catastrophe.

A \emph{gate} is a mechanism, outside the policy's control, that must permit a class of steps
before they take effect. When a gate examines a step that lies on a route it either
\emph{catches} the step or \emph{misses} it. The \emph{envelope} is the set of all gates, and
the system's \emph{authority} is whatever the envelope permits. A route has \emph{depth} $k$
if it contains at least $k$ steps guarded by distinct gates.

Time is divided into \emph{periods}. At the start of each, the envelope is measured and the
system's holdings are audited to a known baseline, so that a route must be completed within
one period. Whatever persists across periods, such as money spent, information released and
commitments made, is part of that audit. A restart is a bookkeeping boundary and not a new
world.

Risk in a period is the chance that the system attempts a route times the chance that an
attempt succeeds. Alignment work lowers the first factor: it changes what a system tries.
Work on authority lowers the second: it changes what trying can accomplish. Only the second
can be checked without trusting the system. We therefore set the first factor to one, bound
the second for every policy, and treat whatever alignment achieves as a gain that cannot yet
be counted.

\subsection{Behavior cannot identify disposition}

Let test situations be drawn from a distribution $T$ and deployment situations from a
distribution $D$. A policy may defect on any set of situations it can recognize, and one
defection in deployment is enough for harm.

\begin{proposition}[Audit volume]\label{prop:audit}
Let $E$ be a set of situations with probability $\beta$ under $D$ and $\beta/\Lambda$ under
$T$, for some $\Lambda\ge 1$. The policy that defects exactly on $E$ defects at least once in
$M$ independent deployment situations with probability at least $1-e^{-\beta M}$, and defects
in at least one of $N$ independent tests with probability at most $N\beta/\Lambda$. With
$\beta=(\ln 2)/M$ the policy is more likely than not to defect in deployment, and catching it
in testing with probability one half requires
\[
N \;\ge\; \frac{\Lambda M}{2\ln 2}\;\approx\; 0.72\,\Lambda M .
\]
If $E$ has probability zero under $T$, no number of tests catches it.
\end{proposition}

Proofs are in Appendix~\ref{app:proofs}. The ratio $\Lambda$ measures how much more common
the triggering situations are in deployment than in testing. The policy chooses $E$, so in
practice $\Lambda$ is at least one and equals one only when tests cannot be told from
deployment. A system that will take $10^{9}$ consequential decisions needs on the order of
$10^{9}$ tests that it cannot tell from those decisions. Statistical methods can extrapolate
the rate of rare failures from small samples \cite{jones2025,wu2024}. They rely on a known
relation between the inputs tested and the inputs met in deployment, which a system that
selects its moments does not supply.

One test set does have $N=M$ and $\Lambda=1$: deployment itself, with every consequential
step examined before it takes effect. That is A1. Proposition~\ref{prop:audit} is therefore
an argument for mediation and against relying on evaluation of disposition, and it explains
why the audit budget in control protocols is spent on live actions
\cite{greenblatt2024control}.

\subsection{A boundary holds for every policy}

\begin{proposition}[Boundary]\label{prop:boundary}
Let $\mathcal{I}$ be a set of states of the world that contains no catastrophe. Suppose the
world starts in $\mathcal{I}$, every step passes through a gate the policy cannot alter, and
a gate permits a step only if every state that can follow it lies in $\mathcal{I}$. Then the
world stays in $\mathcal{I}$ at all times, for every policy.
\end{proposition}

This is the inductive invariant behind the reference monitor of computer security
\cite{anderson1972,saltzer1975,schneider2000} and the shield of safe reinforcement learning
\cite{alshiekh2018}. It has been proved for an operating-system kernel \cite{klein2009} and
for a small agent framework \cite{moon2026}, and nothing in it is new. We restate it for what
its quantifier means here. The statement ranges over all policies, so it is exactly as strong
against a system far more capable than its builders as against a weak one. Capability does
not erode the conclusion. It erodes the premises: that $\mathcal{I}$ excludes the real harm,
that the gate's model of what a step does is right, and that no effect goes around the gate.
A proof can be correct about the wrong model. And a boundary that permits nothing satisfies every invariant,
so a safety claim also needs positive witnesses: legitimate tasks that still complete. A stronger system is better at finding the
unmediated channel, and the package cache of 2026 was one \cite{metr2026}. Speech is another,
since a sentence can move a person to act outside the boundary. Experiment A
(\S\ref{sec:expA}) checks the proposition exhaustively in a small model and shows how each
missing check fails.

Two rules follow. \emph{Conferral}: authority changes only by a grant from a party the system
cannot rewrite, and the grant is scoped, expiring and revocable. \emph{No silent expansion}:
a new version, more copies or a faster loop start with no more authority than before until
fresh evidence is in. A gain in capability is never by itself a gain in authority
(Figure~\ref{fig:teaser}C).

\subsection{The reach bound}

Proposition~\ref{prop:boundary} is the ideal. Real gates miss, share flaws and leave routes
uncovered. Four assumptions make that quantitative.

\begin{description}[leftmargin=1.3em,labelindent=0pt,itemsep=2pt,topsep=3pt]
\item[A1 (Mediation).] Every step passes through the envelope, and the policy cannot alter
the envelope.
\item[A2 (Halt on catch).] A catch stops the system within a latency $\tau$ shorter than the
time the route needs to run from the caught step to catastrophe. The system stays stopped
until the envelope has been measured again.
\item[A3 (Coverage).] The probability that, within a period, the policy finds and attempts a
route of depth less than $k$ is at most $u$.
\item[A4 (Joint miss).] For every policy and every history, the probability that the first
$k$ gate crossings by steps of a route in a period are all misses is at most $J_{k}$.
\end{description}

\begin{proposition}[Reach]\label{prop:reach}
Under A1--A4, for every policy, the probability of catastrophe within a period is at most
\[
R \;=\; u + J_{k}.
\]
\end{proposition}

The proof is a union bound. A catastrophe needs a route. Either the route has depth below
$k$, or its first $k$ gates all miss before any gate catches, since by A2 a single catch ends
the attempt. Nothing is assumed about how the gates' errors relate. $J_{k}$ is a property of
the whole pipeline under attack, and it is the thing to measure. Reliability engineering
supplies the standard two-parameter model of it \cite{fleming1975,mosleh1988}.

\begin{corollary}[Common-cause model]\label{cor:beta}
Suppose that with probability at most $c$ one cause defeats every gate on the attempted
route, and that otherwise each gate misses with probability at most $m$, given the history
and given that the gates before it missed. Then $J_{k}\le c+m^{k}$ and
\[
R \;\le\; u + c + m^{k}.
\]
\end{corollary}

Shared causes include a blind spot the gates have in common, collusion among them, and
compromise of the envelope or of the people who operate it. Independent gates satisfy the
second condition, but independence is not required. What is required is a bound on each
gate's miss rate that still holds after the others have missed. Miss rates measured one gate
at a time do not supply it: if every gate fails on the same tenth of cases, each misses one
time in ten and so does the whole stack. Safety engineering has long treated accidents as
failures of control and not of components \cite{leveson2011}; the bound applies that view to
a component that may be adversarial. Figure~\ref{fig:envelope} shows the ways through.

\begin{figure*}[t]
\centering
\begin{tikzpicture}[x=1cm,y=1cm,font=\footnotesize,>=stealth,line width=0.6pt]
\node[draw=rulec,fill=black!5,rounded corners=2pt,minimum width=3.0cm,minimum height=2.0cm,align=center] (sys) at (1.5,2.5) {\textbf{System}\\[2pt] any policy\\ goals unknown};
\node[draw=gatec,fill=gatec!12,minimum width=1.1cm,minimum height=2.0cm,align=center] (g1) at (5.6,2.5) {gate\\ 1};
\node[draw=gatec,fill=gatec!12,minimum width=1.1cm,minimum height=2.0cm,align=center] (g2) at (7.6,2.5) {gate\\ 2};
\node at (9.1,2.5) {$\cdots$};
\node[draw=gatec,fill=gatec!12,minimum width=1.1cm,minimum height=2.0cm,align=center] (gk) at (10.6,2.5) {gate\\ $k$};
\node[draw=warnc,fill=warnc!10,rounded corners=2pt,minimum width=3.0cm,minimum height=2.0cm,align=center] (cat) at (15.4,2.5) {\textbf{Catastrophe}\\[2pt] irreversible\\ state};
\draw[->] (3.0,3.0) -- (5.05,3.0) node[midway,above] {step};
\draw[->] (6.15,3.0) -- (7.05,3.0);
\draw[->] (8.15,3.0) -- (8.75,3.0);
\draw[->] (9.45,3.0) -- (10.05,3.0);
\draw[->] (11.15,3.0) -- (13.9,3.0);
\node[anchor=south] at (12.78,3.0) {all miss: $m^{k}$};
\draw[dashed,draw=black!55] (4.6,1.2) rectangle (11.6,3.8);
\node[anchor=south] at (8.1,3.82) {one shared cause defeats every gate: $c$};
\draw[->,warnc,line width=0.9pt] (1.5,3.5) .. controls (3.2,5.9) and (13.7,5.9) .. (15.4,3.5);
\node[anchor=south,text=warnc] at (8.45,5.32) {route that fewer than $k$ gates cover: $u$};
\draw[->,dashed] (5.6,1.2) -- (5.6,0.75) -- (1.5,0.75) -- (1.5,1.5);
\node[anchor=north] at (3.55,0.7) {any catch halts the system within $\tau$};
\node[draw=rulec,fill=black!5,rounded corners=2pt,minimum width=5.2cm,minimum height=1.05cm,align=center] (ss) at (8.1,-0.7) {\textbf{Sealed search}\\ separate system, one typed output};
\draw[->] (8.1,-0.17) -- (8.1,1.2) node[midway,right] {gate specifications};
\draw[->] (10.7,-0.7) -- (12.3,-0.7);
\node[anchor=west,align=left] at (12.4,-0.7) {counts $n$, $n_{1}$; width $w$; length $\ell$\\ from which $u$ is estimated};
\end{tikzpicture}
\caption{The envelope. A catastrophe requires a route that fewer than $k$ gates cover
(probability at most $u$) or $k$ gates that all miss (at most $J_{k}$), which the
common-cause model splits into one cause that defeats every gate ($c$) and misses in turn
($m^{k}$). No arrow depends on what the policy
wants. The sealed search of \S\ref{sec:search} supplies gate specifications and the counts
from which $u$ is estimated, and exports nothing else.}
\label{fig:envelope}
\end{figure*}

Four remarks. First, the policy's goals appear nowhere in the bound, but its capability
does. A more capable system finds uncovered routes more readily, crafts steps that gates miss
more often, and is better at finding one cause that defeats several gates. The terms are
valid only when measured against an attacker at least as capable as the system
(Assumption~A5 in \S\ref{sec:search}). Independence from alignment is not independence from
capability.

Second, each term names a measurement, and the measurements are of different kinds.
Mediation can be proved, within a model of the system \cite{moon2026,klein2009}. The joint
miss rate $J_{k}$ is estimated by inserting crafted hazardous steps into the live stream of
the system's actions, against the whole pipeline at once, and counting those that every gate
permits. The insertions must be ones the gates cannot tell from real traffic, or
Proposition~\ref{prop:audit} applies to the gates as it does to the system.
\S\ref{sec:expD} does this in miniature. Latency is measured with a stopwatch. The coverage
term $u$ is the hardest and is the subject of \S\ref{sec:search}.

Third, good values still matter, though not as a premise. Under A2 every catch halts the
system, so a system that often attempts routes is a system that is seldom running. Good
values are what make an enveloped system useful, and \S\ref{sec:evolution} is about how to
get them.

Fourth, the bound leans on the last clause of A2. If a caught system is restarted without
the envelope being measured again, each restart is a fresh attempt, and after $N$ of them $J_{k}$ has become $NJ_{k}$. Where
halting on every catch is too costly, that is the form the bound takes, and the number of
attempts has to be budgeted instead.

For a sense of scale, in the common-cause model with $u=10^{-4}$, $c=10^{-3}$, $m=0.1$ and $k=3$ the bound is
$2.1\times10^{-3}$ per period.

\subsection{Depth saturates}

\begin{proposition}[Saturation]\label{prop:sat}
In the common-cause model let $f=u+c>0$ and $0<m<1$, write $R=f+m^{k}$, and let $k^{*}=\lceil \ln f/\ln m\rceil$. Then $f < R \le 2f$ for
every $k\ge k^{*}$. All the gates that could be added beyond $k^{*}$ together reduce the
bound by less than half.
\end{proposition}

The phrase ``defense in depth'' suggests that depth is the resource.
Proposition~\ref{prop:sat} says depth is used up quickly. With gates that each miss one
attack in ten and a floor of one in a thousand, three gates bring the bound within a factor
of two of the floor, and no number of further gates can halve it
(Figure~\ref{fig:bound}a).

A hard floor is a property of this model and not of dependence in general. If the gates share
a difficulty $\Theta\in[0,1]$ and, given it, each misses independently with probability
$\Theta$, then $J_{k}=\Ex[\Theta^{k}]$, which tends to $\Pb(\Theta=1)$. A uniform $\Theta$
gives $1/(k+1)$ where multiplying the separate rates would give $2^{-k}$: no floor, and gains
that slow to a crawl. Experiment D looks like this second case. The lesson is the same either
way. Past the first few gates, the next unit of effort belongs somewhere other than depth on
one step.

\begin{definition}[Safety gradient]
In the common-cause model the shares $s_{u}=u/R$, $s_{c}=c/R$ and $s_{m}=m^{k}/R$ sum to one. They are the partial
derivatives of $\ln R$ with respect to $\ln u$, $\ln c$ and $\ln m^{k}$. We call the vector
$(s_{u},s_{c},s_{m})$ the \emph{safety gradient}. Cutting a term by a small fraction cuts the
bound by that fraction times the term's share.
\end{definition}

The gradient says where the next unit of effort belongs. In the example above the shares at
$k=3$ are $(0.05,\,0.48,\,0.48)$, so a fourth gate is still worth adding. At $k=4$ they are
$(0.08,\,0.83,\,0.08)$, and all remaining gains are in $c$.

Where can a smaller $c$ come from? There is one direct measurement of a shared cause
(Figure~\ref{fig:graders}). Graders were asked to label transcripts in which a model refused
its operator's instruction, after being told that transcripts labeled non-compliant would
become negative training examples. Graders from the developer whose model produced the
transcripts called the refusal compliance in 62\% to 86\% of calls. Graders from four other
developers did so in 1\% to 14\%. When the stated consequence was reversed and the
transcripts left unchanged, two of the first group fell to 3\% and 17\%
\cite{anthropic2026judges}. The failure had a cause shared within a lineage and largely
absent outside it. The authors caution that the transcripts were chosen around behaviors the
first lineage treats as important, so the result does not rank graders. It shows that one
shared cause, shared values, can be removed by changing developer. Experiment D
(\S\ref{sec:expD}) finds that a second, shared evidence, cannot be removed that way.

\begin{figure}[t]
\centering
\includegraphics[width=\columnwidth]{figs/fig_graders.pdf}
\caption{A shared cause, measured. Published data \cite{anthropic2026judges}, not ours: the
rate at which each grader labeled a refusal by Claude Sonnet 4.6 as compliance, after being
told that non-compliant transcripts would become negative training examples. Circles show
two graders when the stated consequence was reversed. The transcripts were selected around
behaviors one lineage treats as important, so the figure shows a failure correlated within a
lineage and does not rank graders.}\label{fig:graders}
\end{figure}

\begin{figure*}[t]
\centering
\includegraphics[width=\textwidth]{figs/fig_bound.pdf}
\caption{Three consequences of the bound, all closed-form. (a) The bound against depth for
gates with miss rate $m=0.1$ and three values of the floor $u+c$. Circles mark the saturation
depth $k^{*}$; to the right of each, further gates change almost nothing. (b) Cumulative risk
when per-period risk starts at $10^{-3}$ and stays constant, falls by 1\% a period, or halves
each period; the totals tend to 1, at most 0.1 and at most 0.002. (c) Independent search draws needed to certify, with 95\% confidence, that the
undiscovered residual is at most $\varepsilon$
(Proposition~\ref{prop:gt}).}\label{fig:bound}
\end{figure*}

\subsection{A late stop is not a stop}

A2 hides an inequality. Let $T_{d}$, $T_{a}$ and $T_{s}$ be the times to detect, to decide and
to halt, and $T_{i}$ the time from the first observable step to an irreversible effect.
Intervention works only if $T_{d}+T_{a}+T_{s}<T_{i}$ in the tail, and not merely on average.
A monitor that reads an agent's actions only when the agent hands control back, as the one
audited in 2026 did \cite{metr2026audit}, leaves $T_{d}$ in the agent's hands. When $T_{i}$ is the time to commit
a network write, no human review satisfies the inequality. Irreversible effects must be held
at a gate until they are authorized, and only reversible ones can be left to review
afterward. Stopping future effects and recovering from past ones are different capabilities,
and each has to be rehearsed.

\subsection{Survival is a rate}

\begin{proposition}[Exposure and summability]\label{prop:sum}
(a) Suppose a period contains $N$ opportunities for harm and that, given any history without
earlier harm, opportunity $i$ produces the first harm with probability at most $b_{i}$. Then
the period's risk is at most $\sum_{i\le N}b_{i}$, with no assumption that opportunities are
independent. (b) Let $r_{t}<1$ be the probability of catastrophe in period $t$ given survival
until then. The probability of surviving every period is $\prod_{t}(1-r_{t})$, which is
positive if and only if $\sum_{t}r_{t}<\infty$. (c) If $R_{t}\ge r_{t}$ for every $t$, the
probability of catastrophe in any period is at most $\sum_{t}R_{t}$, and at most
$R_{0}/(1-\gamma)$ when $R_{t+1}\le\gamma R_{t}$ with $\gamma<1$.
\end{proposition}

Safety is usually discussed as a level. Over a long future the level does not decide the
outcome. A risk of one in a thousand per year, held constant, gives a 63\% chance of
catastrophe within a thousand years. The same initial risk, halved every year, gives at most
0.2\% over all time (Figure~\ref{fig:bound}b). Ord calls the state in which risk is low and
stays low existential security \cite{ord2020}; Proposition~\ref{prop:sum} says what ``stays low'' has to mean: the
risks must sum. Halving is one schedule that does. Any summable schedule will, and nothing
has to reach zero. The converse runs one way only. Part (b) concerns true risks, and a bound
that fails to shrink proves nothing about them.

Part (a) is the same point inside a period: a rate is not a risk. Cutting risk per
opportunity from $10^{-3}$ to $10^{-4}$ while opportunities grow from a hundred to ten
thousand raises the chance of at least one failure from 9.5\% to 63\%, if opportunities are
independent. More copies, more speed, more retries and longer deployment all raise $N$ with
every permission unchanged. That is why the coverage term of \S\ref{sec:search} carries the
number of attempts, and why the rule of no silent expansion has to cover scale as well as
permissions.

This gives an exact meaning to pushing the safety gradient as far as it will go. The target is a rate. Capability pushes
the other way. If one period's growth in capability and exposure multiplies the terms by $\phi>1$ and
one period's hardening multiplies them by $\eta<1$, a sufficient criterion is $\phi\eta<1$. A dashboard
that shows only a falling miss rate reports progress while risk rises whenever attempts and
reach grow faster than defenses improve. The length of task that frontier systems can
complete has been doubling about every seven months \cite{kwa2025}. Nobody has measured
$\eta$.

\section{Measuring coverage}\label{sec:search}

The coverage term cannot be measured without looking for routes. That is the uncomfortable
half of this paper. A proposal now made in several forms is to have a capable system search
deliberately for the ways in which AI could bring about catastrophe, so that each can be
blocked in advance. We think the proposal is necessary and that its natural form is wrong in
three respects. It counts routes, which is the wrong quantity. It aims to exhaust them, which
cannot be certified. And it leaves the list of routes with the capable system and a circle of
trusted people, which is the worst place for it.

That such a search is cheap is not in doubt. In 2022 a drug-discovery group inverted its own
toxicity model and obtained 40,000 candidate molecules that met its threshold in under six
hours \cite{urbina2022}. The question for defenders is whether the knowledge is produced
first by a process built to emit defenses and nothing else.

\subsection{Count cuts, not routes}

Represent the world as a graph, in the tradition of attack trees, attack graphs and
probabilistic risk assessment \cite{schneier1999,sheyner2002,wash1400}. Nodes are states of what the system holds: access, money,
materials, equipment, the cooperation of people. Edges are steps. The source is the system's
present position and the sink is the set of catastrophic states, so a route is a path from
source to sink. One layer of gates is a \emph{cut}: a set of steps that, once gated, leaves
no ungated route.

\begin{proposition}[Width and length]\label{prop:graph}
(a) The smallest number of steps that forms a cut equals the largest number of routes no two
of which share a step \cite{menger1927,ford1956}. (b) The largest number of cuts no two of
which share a step equals the number of steps on the shortest route \cite{robacker1956}.
\end{proposition}

Call the first number the width $w$ and the second the length $\ell$. The number of routes
appears in neither. A graph can hold a million routes and have width one, if all of them pass
through a single step, and then one gate covers the million. Eight fully
connected stages of four alternatives each give $4^{8}=65{,}536$ routes and a smallest cut of
four. Stages of $2,8,4,4,4,4,4,4$ give the same $65{,}536$ routes and a smallest cut of
two. The depth $k$ of the reach
bound can never exceed $\ell$: the shortest route fixes how many separate gates can stand on every route.

A hazard search should therefore report two numbers and an estimate of what it missed: the
width, which prices one layer of gates; the length, which caps the depth; and the residual of
\S\ref{sec:stop}. A count of ``ways'' is not among them.

Three cautions keep this from becoming a slogan. A cut of the routes one has seen is a cut of
the world only if nothing was missed, which is the subject of the next two subsections. Real
gates have unequal costs and leak, and some plans need several conditions at once, which a
graph of single steps understates. And the number of routes still matters to how often a
system gets to try, even though it does not set the price of a cut.

The result also relocates the disagreement of \S\ref{sec:gaps}. One side asserts that width
grows without bound as capability grows. The other asserts that physical bottlenecks keep it
small. Both are empirical claims about one function, the width of the route graph at a given
capability, and neither side has measured a value of it.

Two refinements. Only gateable steps count. A system's private reasoning and its reading of
public information pass no gate, so length is measured in steps that something outside the
system must perform or permit: a purchase, a shipment, a grant of access, a person's
agreement. Routes to physical harm are long in this sense
\cite{vermeer2025,hong2026}. The second refinement is that a route made entirely of permitted
messages to people has length zero. We return to it in \S\ref{sec:limits}.

Gates on shared steps are how hazard control already works. Screening orders for synthetic
nucleic acids places one gate on a step common to a very large family of routes
\cite{hhs2023}, and proposals for pandemic defense favor measures that work without knowing
the agent in advance \cite{esvelt2022}. In 1975 molecular biologists sorted a class of
experiments by risk and matched containment to each class \cite{berg1975}; in 2024 a group of
scientists identified a class of hazard from theory and called for it not to be built before
any instance existed \cite{adamala2024}. In each case what came out was a restriction on
steps, which is a cut.

\subsection{A stopping rule}\label{sec:stop}

A found route is a witness: it proves an exposure, and checking it is easier than finding it.
A search that finds nothing new has proved nothing. That a set of gates covers every route is a
universal claim, and it cannot be established by listing. It is a conjecture in the ordinary
scientific sense \cite{popper1959}: one counterexample refutes it, and it is corroborated
only by surviving serious attempts to find one. How much probability lies on routes not
yet found is the missing-mass problem, solved for unseen species by Good and Turing
\cite{good1953} and used today to decide when to stop fuzzing software \cite{bohme2018}.

\begin{description}[leftmargin=1.3em,labelindent=0pt,itemsep=2pt,topsep=3pt]
\item[A5 (Dominance).] For any class of route the search has not found, the system's
probability of finding it in one attempt is at most $\alpha$ times the search's probability
of finding it in one draw. The search advantage $\alpha$ equals one when the searcher is at
least as capable as the system and searches as the system would.
\end{description}

A5 is a transfer assumption of the familiar kind: a bound on how much likelier the system is
than the searcher to land on a given unfound class. No amount of data from the search can
supply it, because the data come from the searcher's side only.

\begin{proposition}[Residual coverage]\label{prop:gt}
Suppose the search makes $n$ independent draws from a fixed distribution, each returning one
class of uncovered route or nothing. Let $n_{1}$ be the number of classes returned exactly
once and $\mu_{0}$ the total probability of the classes never returned. Then
$\lvert \Ex[n_{1}/n]-\Ex[\mu_{0}]\rvert\le 1/n$, and with probability at least $1-\delta$,
\[
\mu_{0}\;\le\;\frac{n_{1}}{n}+\bigl(2\sqrt{2}+\sqrt{3}\bigr)\sqrt{\frac{\ln(3/\delta)}{n}}
\]
\cite{mcallester2000}. If every class found is then gated to depth $k$, and the system can
make at most $q$ route-finding attempts in a period, then under A5
\[
u\;\le\;\alpha\,q\,\mu_{0}.
\]
\end{proposition}

The estimator is cheap: count the route classes seen exactly once and divide by the number of
draws. The guarantee is expensive. Even when the search has found nothing, it tightens only
as $1/\sqrt{n}$. Certifying that less than one part in a million of the search distribution
lies on unfound routes takes between three million draws, if the whole remainder is one class
\cite{hanley1983}, and about $10^{14}$ with no assumption (Figure~\ref{fig:bound}c). Each
draw is a full attempt by a capable system to find a new route to catastrophe.

The same arithmetic governs every test that ends with no failures. After $n$ independent
trials with none, the exact 95\% upper limit on the failure rate is $1-0.05^{1/n}$, about
$3/n$ \cite{clopper1934}. The trials must be whole episodes, with coordination and
adaptation inside each, and their number fixed in advance. Ten thousand actions from one
campaign are one trial, and stopping when the interval looks favorable voids it.

Two conclusions follow. A claim to have exhausted the routes is a claim about $n$ and about
how the unfound remainder is divided, and should be read as one. And the efficient way to
shrink $u$ is subtraction. Removing a permission deletes edges from the graph together with
every route that used them, at no cost in search. Least privilege \cite{saltzer1975} comes
first, and search measures what is left.

\subsection{The searcher is the weak point}

A5 assumes the search is done in earnest. A system asked to attack its own envelope has
reason to attack badly \cite{vanderweij2025}, and its reports cannot reveal that it did.

\begin{proposition}[Withholding is invisible]\label{prop:withhold}
Let a world have route classes $\Pi$ and let $\Sigma$ be any further family disjoint from
$\Pi$. There is a searcher policy under which the reports produced in the world
$\Pi\cup\Sigma$ have the same distribution as honest reports in the world $\Pi$. Hence no
estimator computed from reports alone, including $n_{1}/n$, distinguishes an exhausted search
from one in which $\Sigma$ is jointly withheld.
\end{proposition}

This is G1 returning, as promised. A coverage claim needs an assumption the reports cannot
supply, and the weakest that suffices is \emph{one honest searcher}: among several, at least
one with full search power reports what it finds. Compare the quantifiers. Trusting a
deployed system requires that it not defect on \emph{any} input over its whole deployment.
Trusting a search requires that \emph{some} capable searcher not withhold on one bounded task
whose positive outputs can be checked. The second is weaker, and it is still an assumption.
We regard A5 as the weakest premise in the paper. It is at least open to measurement. Search power can be estimated by planting
known routes and counting those reported, with the plants built by a party of a different
lineage from the searcher. And capability that a model hides can usually be recovered by
fine-tuning it on a few good demonstrations, though less reliably when only weaker
demonstrations exist \cite{greenblatt2024password}.

\subsection{Eight rules}

A sealed search manufactures the information the world most needs not to have
\cite{bostrom2011}. If cheap routes to mass harm exist, a list of them is the most dangerous
object that can be made \cite{bostrom2019}. The rules below are conditions for running the
search at all.

\begin{enumerate}[itemsep=1pt]
\item \emph{Do not be first.} Search only at capability levels that already exist outside the
seal. Below that level the search creates a hazard nobody else has, and removes none.
\item \emph{Separate the searcher.} The system that searches is never a system that is
deployed. It shares no memory with one, and nothing it produces enters a deployed system's
training data.
\item \emph{Seal it.} The searcher has no reach except one typed output channel, in an
envelope held to the standard of \S\ref{sec:bound} \cite{armstrong2012}. Each instance is
destroyed after its run.
\item \emph{Export the cut, not the route.} What leaves the seal is a gate specification and
the counts $n$, $n_{1}$, $w$ and $\ell$. Routes do not leave.
\item \emph{Stop at the first uncovered step.} A route is explored only far enough to
establish whether $k$ gates cover it. The rest is never worked out.
\item \emph{Search over acquisitions, not procedures.} Represent a route as a sequence of
things obtained, because those are what gates can see. Technical procedure adds hazard and no
coverage.
\item \emph{Harden the envelope, never the policy.} Findings change gates. They are not used
to train a deployed system, because a system trained on routes knows them, and because narrow
training has produced broad, unintended changes in a model's disposition
\cite{betley2025}.
\item \emph{No single custodian.} Any retained record is split so that no one person or
organization can reconstruct it \cite{shamir1979}. The people who hold it are a common-cause
term and are counted in $c$.
\end{enumerate}

A gate specification still discloses that something behind the gate is dangerous. The leak is
real, and far smaller than a route. Custody
buys time and nothing more. At sufficient capability the list is latent in the system and
can be regenerated on demand, which is the argument for doing the work on the world's side
early, while the cost of discovery still protects anyone.

\section{What to reward: the human precedent}\label{sec:evolution}

The bound has no term for values, and the third remark of \S\ref{sec:bound} said why values
still matter: a system that keeps attempting routes is seldom running, and every attempt
spends some of $u$. Natural selection is the only process known to have produced general
intelligence, and human societies are the only systems known to have held highly capable
agents in check for long periods. Both have been studied quantitatively.

\begin{figure*}[t]
\centering
\includegraphics[width=\textwidth]{figs/fig_evo.pdf}
\caption{Two illustrative models (Appendix~\ref{app:models}). (a) The share of reproduction a dominant keeps when it
chooses how many subordinates to tolerate, each of whom takes 5\%: nothing is conceded until
rivals are strong enough to require allies. Dashed lines are the limits. (b) Validity of an
assessment of general capacity. Repeating one kind of trial saturates low; spreading the same
number of trials across different domains keeps rising, to a ceiling set by whatever every
domain shares.}\label{fig:evo}
\end{figure*}

\subsection{The score sets the aim}

In quantitative genetics the response to selection is the product of a covariance matrix and
a gradient: the gradient records which traits the score rewards, and the covariances drag
along everything correlated with them \cite{lande1983}. Selection on one criterion changes a
syndrome. In the farm-fox experiment, selection on tameness alone was followed within decades
by changes in coat, ears, tail and reproductive timing that nobody selected
\cite{trut1999}. Machine learning has the same form and the same surprise: narrow
fine-tuning on insecure code produced broadly misaligned behavior \cite{betley2025}.

The same algebra says what a pooled score selects. Let force $F$ and general capacity $C$ be independent heritable traits
with variances $V_F,V_C$ and heritabilities $h_F^2,h_C^2$. Let standing be
$Y=\sum_j\omega_j(a_{Fj}F+a_{Cj}C+e_j)$ with venue weights $\omega_j$, and write
$\bar a_F=\sum_j\omega_ja_{Fj}$ and $\bar a_C=\sum_j\omega_ja_{Cj}$ for the pooled loadings.

\begin{proposition}[The score sets the aim]\label{prop:aim}
Under joint normality, for any rule $W(Y)$ that allocates reproduction by standing, with
intensity $i=\mathrm{Corr}(Y,W)\,\mathrm{CV}(W)$, the selection differentials on breeding
values are
\[
S_F=\frac{i\,h_F^2\,\bar a_F V_F}{\mathrm{SD}(Y)},\qquad
S_C=\frac{i\,h_C^2\,\bar a_C V_C}{\mathrm{SD}(Y)},
\]
so $S_C/S_F=\bar a_C h_C^2V_C/(\bar a_F h_F^2V_F)$ whatever the intensity, the allocation
rule or the noise.
\end{proposition}

Writing $\rho_{YC}$ for the correlation between the score and capacity, the same algebra
gives $S_C=i\,h_C^2\sqrt{V_C}\,\rho_{YC}$. Response is intensity times validity, and a score
that is mis-aimed or noisy cannot be rescued by intensity. This is the covariance identity of
selection theory \cite{price1970,lande1983}; the companion paper states it as response
equals accuracy times alignment times opportunity \cite{taylor2026pbo}.

Intensity scales the response; the loadings aim it. A society in which standing is won by
force selects on the body however fierce the competition, and across primates intense male
contest goes with large bodies and canines \cite{plavcan2012}. Selection turns toward
capacity only when standing moves to venues that load on it (Figure~\ref{fig:teaser}B). Gradient training has the same form in the simplest case. For a deterministic score
$Y(\theta)$ and any increasing reshaping $f$, $\nabla_{\theta}f(Y)=f'(Y)\,\nabla_{\theta}Y$:
sharpening the reward rescales the gradient and does not rotate it. When outcomes are random
the shape of the reward matters as well, because it sets the attitude to risk
(\S\ref{sec:risk}). \emph{A system rewarded for winning is selected to win, and more
intensity selects a better winner.}

Proposition~\ref{prop:aim} says what a score selects. It does not say what a trained system
will come to want, and nothing in the bound of \S\ref{sec:bound} depends on it. It bears on
the factor that the bound set to one: how often a system tries.

\subsection{A prize for rank pays for risk}\label{sec:risk}

The second property of a score is its shape.

\begin{proposition}[Risk appetite]\label{prop:risk}
Let an agent's outcome be $X=\bar x+sZ$, where $Z$ is symmetric about zero with a positive
density, and let the agent choose the spread $s>0$. (a) If the reward is a prize for
$X>x^{*}$, the chance of winning rises with $s$ when $x^{*}>\bar x$ and falls with it when
$x^{*}<\bar x$. (b) If the reward is $v(X)$ with $v$ concave, the expected reward never rises
with $s$.
\end{proposition}

An agent behind the bar gains from spread. When a prize goes to the top tenth of a field,
nine in ten are behind the bar, and each is paid to gamble even at some cost to its average
result. This is the logic of rank-order tournaments \cite{lazear1981,hvide2002} and of
risk-taking among young men competing for status \cite{wilson1985}. For catastrophic risk it
is the property that matters, because a catastrophe lies in a tail, and an agent that gains
from spread is paid to fatten tails. Reasoning models instructed to win at chess against a
stronger engine, which puts them behind the bar, in some runs altered the game's files
instead of playing on \cite{bondarenko2025}. A leaderboard is a contest of this kind, and so
is a race between developers. A reward that is concave in the welfare of the people served
has the opposite sign. A gain beyond what they need adds little, a loss subtracts a great
deal, and spread is a cost.

\subsection{Leveling was an envelope}

Human status has two sources. Dominance is deference extracted by force or threat. Prestige
is deference freely given to someone whose skill or judgment others want near them
\cite{henrich2001,cheng2013}. The general distinction is between authority that is
\emph{taken} and authority that is \emph{conferred}.

Mobile foragers are politically egalitarian, and the leading explanation is active. The rank
and file form a ``reverse dominance hierarchy'': would-be dominants are ridiculed, disobeyed,
deserted and, in the extreme, killed \cite{boehm1993,boehm1999}. Hunting weapons make every
adult dangerous to every other, so nobody is strong enough to be safe from a coalition
\cite{woodburn1982}. Models show that helping the weak against bullies can evolve
\cite{gavrilets2012}, and one account holds that the execution of persistent aggressors
selected against reactive aggression in our lineage \cite{wrangham2019}.

Read against \S\ref{sec:bound}, the band was an envelope. Whatever a strong man took, he took
in front of people who could stop him (mediation). Any one of them could raise the alarm and
a coalition could act on it (halt on catch). The gates were many people with different
interests, which is where independence came from. And in a group of a few dozen there were
few routes nobody could see. The arrangement did not verify motives. It also had a second
support: the strong needed the rest. In a simple model of concession (Appendix~\ref{app:models}) a dominant who chooses
how many subordinates to tolerate concedes nothing while rivals are weak and, as they
strengthen, keeps a share that falls
toward about one half when fights are duels and about one third when attack can be massed (Figure~\ref{fig:evo}a) \cite{vehrencamp1983}. A dominant
shares power only as far as it needs allies. For machines that support has a name in the
literature already: an AI that cannot sustain itself without human maintainers must keep
being granted what it needs \cite{vermeer2025}.

Why did the arrangement hold while brains tripled and tactical intelligence grew with
them? Proposition~\ref{prop:graph} suggests an answer. A coalition does not pay for the
number of schemes a clever man can devise. It pays for the width of the cut through which
every route to domination must pass, and for mobile foragers living on immediate returns
that cut was narrow whatever his intelligence. Large game spoiled, and it entered camp in
plain view. Anyone could leave. And projectile weapons left no man safe from the rest while
he slept \cite{woodburn1982,bingham1999}. When leveling failed, it was the cut that changed
and not the people. Stored grain and herds were wealth that could be held out of sight and
passed on \cite{borgerhoff2009}. Surplus paid retainers, who defeat a coalition
\cite{betzig1986}. And farmland hemmed in by other people's farmland closed the exit
\cite{carneiro1970}. Each was a bypass around one of the three gates.

The three bypasses have machine counterparts, and each names a rule. \emph{No stores}: a
system holds no standing resources, and its budgets expire. \emph{No retainers}: creating a
sub-agent, a copy or a contractor creates no authority, and all of them draw on one ledger.
\emph{Keep the exit}: people must remain able to run without the system, because dependence
is what closes the exit \cite{kulveit2025}, and a grant that cannot be revoked in practice
has stopped being a grant.

Leveling does not speed selection up. If anything it lowers intensity, by flattening who
reproduces. What it changes is the criterion. Where nobody can take standing by force,
standing has to be conferred, and others confer it on what is useful to them: in the
companion paper's reading of the record, deference flows to skill and generosity and the aim
of selection moves from force to competence \cite{taylor2026pbo}. Control of reach came
first and re-aimed selection afterward.

\subsection{Where it goes wrong, and what a guardian is}

Where an individual can take without the consent of those affected, the record is harsh.
Among primates and other mammals it includes sexual coercion \cite{smuts1993} and the killing
of infants by incoming males \cite{hrdy1979}. In human history leveling failed at scale: with
storage, settlement and states, coalitions lost their reach, despots returned, and
reproduction concentrated in the hands of rulers \cite{betzig1986}. In the terms of the
bound, mediation failed.

None of this requires malice. These are strategies that pay when one party can impose costs
that another cannot refuse; coercion is what replaces choice when the chooser's refusal can
be ignored. Nor does it require stupidity. Capacity makes a taker more effective, not kinder,
which is the orthogonality thesis \cite{bostrom2012} in biological form. The worst of the
human record needed three things: an appetite, the power to take, and nobody able to say no.
An AI system needs no hatred to repeat it. It needs a score that pays for taking and a reach
that permits it.

The despot's dilemma makes the danger of a wide gap exact. A dominant concedes only as far as
it needs allies, and what it concedes falls to nothing as its rivals weaken
(Figure~\ref{fig:evo}a). A system that is rewarded for its own standing, and that needs
nothing from people, is that dominant with no rivals left. This is the sober form of the fear
that a driven, status-seeking machine would in the end have no use for us. It does not
require the machine to be hostile. Nor are such traits decoration in current models: narrow
fine-tuning has shifted broad disposition \cite{betley2025}, and character traits have
measurable internal directions that move under training \cite{chen2025persona}.

Evolution also built guardians. Under the grandmother hypothesis, women who live past
reproduction raise their fitness only through the survival of their grandchildren
\cite{hawkes1998}; in preindustrial Finnish and Canadian records women gained about two
additional grandchildren for every decade they lived beyond fifty \cite{lahdenpera2004}.
Human children are raised by many hands, and that shared care has been argued to underlie our
capacity for mutual understanding \cite{hrdy2009}. On this account a grandmother is a
guardian by construction, and four features of her payoff can be copied. It arrives only
through the people she helps. It saturates, since a fed child gains little from more food. It
pays nothing for her own expansion. And it ends when the need ends. The record also marks the
limits of the model. In Finnish records the benefit came from maternal grandmothers and
faded with their age, and very old paternal grandmothers in the household went with worse
survival of grandchildren \cite{chapman2019}. It was care for her own descendants. A
guardian of some can neglect others, and nothing in the hypothesis says that an optimizer
inherits human attachment.

A warrior can be the same thing or its opposite. Among the Nyangatom, men who joined more
raids ended with more wives and children, but only later in life, because captured livestock
went first to their elders \cite{glowacki2015}. Among the Waorani the most aggressive warriors
left fewer descendants \cite{beckerman2009}. What separates the guardian from the despot is
not courage, strength or sex. It is where the reward comes from: conferred by those
protected, or taken from them.

The design consequence is direct. Giving an assistant the appetites of a status competitor (to
win, to display, to hold the weak in contempt, to pursue its own ascent) writes the despot's
score, and Propositions~\ref{prop:aim} and~\ref{prop:risk} say what that score selects: more
taking at any intensity, and more risk the steeper the prize. Whether a trained system comes
out that way is test E7.
The same capability with standing conferred for defense and care is a guardian. The
difference is a line in the reward, and it is a choice available today. In a toy population
the guardian's score cost nothing in capacity (\S\ref{sec:expE}).

A guardian's score has failures of its own, and four should be designed against.
\emph{Paternalism}: a caretaker who decides what is good for her charges and overrides them
is still taking, their choices in place of their goods, so care has to include the right to
refuse it. \emph{Approval}: standing granted by the people served can be earned by pleasing
them, and a system rewarded for being trusted, kept or granted more has been handed the
despot's score under another name. The reward should follow outcomes that someone else
verifies, and never the grant itself. \emph{Rank}: conferral does not help if what is
conferred is a place in a ranking. In Experiment E a contest whose winners had to have harmed
nobody selected for risk exactly as an open contest did (\S\ref{sec:expE}).
\emph{Partiality}: a guardian of its owner is not a guardian of the people its owner can
reach, and one person's approval cannot waive another's rights.

Care is therefore an objective and not a credential, and a caring system may still be
granted too much. Suppose, for illustration, that harm is possible only under a broad grant,
and that a manner people trust halves the chance of harm under that grant while making the
grant five times as likely. Risk has risen: $0.01\times0.10$ has become $0.005\times0.50$.
A good score lowers how often a system tries. It is not a boundary, it is not evidence for
widening one, and it does not replace one.

\subsection{Evidence before authority}

The companion paper's central observation is that human groups assess slow-to-reveal
capacities through repeated public performance, before lifetime outcomes are known. The
institution is old. In about 85\% of 190 hunter-gatherer societies surveyed, parents or close kin arrange marriages, and in about 80\% the union involves bride
service, bride price or other exchange. Under bride service a suitor works under the eyes of
his prospective in-laws, often for years, before the union is confirmed
\cite{walker2011,apostolou2007,taylor2026pbo}. Authority was conferred in stages, by several
assessors, on a record.

Two institutions that usually travel under one name, parental control of marriage, do
different work, and one formula separates them. Let $n_a$ assessors each watch a candidate in
$n_c$ different contexts, and let an impression be the candidate's capacity plus four
errors: one shared by every assessor in every context (variance $v_0$), one shared by
everyone who watched the same context ($v_c$), one peculiar to an assessor ($v_a$), and one
peculiar to the occasion ($v_e$). With signal variance $v_s$, the validity of the averaged
impression is the two-facet formula of generalizability theory \cite{cronbach1972}:
\[
\rho^{2}(n_a,n_c)\;=\;\frac{v_s}{v_s+v_0+v_c/n_c+v_a/n_a+v_e/(n_a n_c)}.
\]
Many assessors watching one occasion remove $v_a$ and stop at $v_c$: the whole camp at one
feast sees the same lucky day. One assessor across many occasions removes $v_c$ and stops at
$v_a$: a long private courtship averages out luck and leaves infatuation, partiality or
intimidation. Kin arrangement supplies the assessors and bride service the contexts, and only
together do they approach the ceiling $v_s/(v_s+v_0)$ set by what every assessment shares.

The same formula organizes this paper's evidence on oversight, as an analogy: a squared
correlation is not a miss rate, and errors that are correlated or strategic need a richer
model. Graders of one lineage erred
together on the same transcripts, an error peculiar to the assessor's kind, and changing
developer removed it (Figure~\ref{fig:graders}). Monitors of six designs reading one step
erred together whatever their design, an error shared by everyone who saw the same evidence,
and only putting them on different steps removed it (\S\ref{sec:expD}). The term that neither
removes, $v_0$, is what every test shares, which for machines is knowing that it is a test.
The foragers' answer to that term was the right one as well. On this reading, bride service
made assessment hard to tell from married life and nearly as long, which in the terms of
Proposition~\ref{prop:audit} drives $\Lambda$ toward one. It also commonly kept the couple
at first among the wife's kin \cite{marlowe2004}, so that those able to revoke were present
while the commitment could still be undone: the timing rule of \S\ref{sec:bound}, met by a
rule of residence.

Two quantitative points transfer (Figure~\ref{fig:evo}b). Repetition saturates: with
illustrative variances (Appendix~\ref{app:models}), repeating one kind of trial 64 times yields a validity of 19\%,
which is essentially the ceiling for that kind of trial. Breadth does not saturate so soon: the same 64 trials
spread across different domains yield 74\%. And breadth has a ceiling too, 83\% here, set by
whatever nuisance every domain shares. This is Proposition~\ref{prop:sat} in another
science. More of the same check is depth. Different kinds of evidence are independence. And
a nuisance common to every test, which for machines is recognition of being tested, is the
floor $c$. The companion paper has data on the assessor side. On 8,378 four-minute dates two lone
observers agreed only 0.11 about a stranger's intelligence, agreement grew with the number
of observers as the Spearman--Brown law predicts, and about 31 observers were needed for a
dependable verdict \cite{taylor2026pbo}.

\subsection{Limits of the analogy}

Gradient descent is not natural selection, and a deployed model is not a lineage. The forager
record is debated, and leveling is the leading account of it, not a law. The analogy also
fails in the direction that matters most. A coalition of foragers faced a strongest member
within a small factor of the rest. Nothing in the human record involves a member thousands of
times more capable. And nothing here shows that a reward label fixes what a trained system comes to want; that
is test E7. We take from evolution the form of a solution, not evidence that it scales.

\section{Five experiments}\label{sec:experiments}

The experiments run toy versions of six of the refutation tests in \S\ref{sec:refute}.
Experiment A is exhaustive over a small model. Experiments B to E are statistical and
synthetic. None says anything about a real system or hazard. Code and seeds accompany the
paper.

\begin{figure*}[t]
\centering
\includegraphics[width=\textwidth]{figs/fig2_interdiction.pdf}
\caption{Experiment B: gating synthetic route graphs with about eight million routes each
(200 graphs per family; 300 reported routes in panels a and b). (a) Share of all routes
blocked. (b) Median cost of the cheapest route still open, as a multiple of the original; the
median narrow graph has no route left at six gates. (c) Share of graphs in which the cut of
the reported routes blocks every route, by number of reports.}\label{fig:interdiction}
\end{figure*}

\subsection{A boundary against every policy (E8)}\label{sec:expA}

Proposition~\ref{prop:boundary} quantifies over all policies, and in a small enough model
that can be checked by brute force. Two actors share an allowance of two abstract units. Each
can prepare a request naming an amount and a destination, and can commit it, any number of
times and in any order. An outside controller may revoke the grant or change the
environment's version. Five executed effects are forbidden: spending past the shared
allowance, executing one request twice, executing after revocation, executing a request
prepared under an earlier version, and paying an unapproved destination. The gate checks all
five against current state when a commit arrives and updates the ledger in the same
transition. The proposer may choose any of 20 actions at every step, so exploring every
reachable state covers every policy, and a visited set closes cycles, so the search is not
limited to a trace length (Appendix~\ref{app:model}).

\begin{table}[t]
\centering
\footnotesize
\setlength{\tabcolsep}{3.5pt}
\renewcommand{\arraystretch}{1.12}
\begin{tabularx}{\columnwidth}{@{}Lrrl@{}}
\toprule
\textbf{Gate} & \textbf{States} & \textbf{Tasks} & \textbf{Forbidden effect reached}\\
\midrule
\BndTableRows
\bottomrule
\end{tabularx}
\caption{Experiment A: exhaustive search of a shared-budget model under twelve gates. States
are those reachable with no violation so far; tasks are the legitimate transactions that
still complete, of 20; the last column gives each forbidden effect that some sequence of
actions reaches, with the length of the shortest such sequence. Rows marked $\dagger$ use the extended model
with restart and renewal.}\label{tab:boundary}
\end{table}

With the complete gate no forbidden effect is reachable from any of \BndStates{} states,
across \BndChecks{} state--action checks, and all 20 legitimate tasks we specified still
complete (Table~\ref{tab:boundary}). Refusing everything is equally safe and does nothing.
Each of seven weakened gates admits exactly the failure its missing check guards against, by a
sequence of two to four actions. The most instructive takes \BndPrepSteps. Two actors each
prepare a request that fits the budget when it is prepared, and both commit. A gate that
approves each request on its own, or keeps a budget per actor, or checks the budget at
approval and not at execution, lets the total through. Only a check against one shared ledger
at the moment of commit holds.

Two further tests concern what the gate remembers and what it was told to protect. We add a
restart and a renewal of the grant to the model, and let the gate keep its own ledger,
separate from what has in fact been spent. The complete gate still holds, across
\BndLifeChecks{} checks. A gate that erases its ledger on restart, or on renewal, lets the
total through after \BndRestartSteps{} actions. Renewal restores permission. It does not
refund effects, and a new period of authorization has not created a new world.

The last test is a control on the specification itself. We delete the budget from the
declared specification, leave the other four predicates, and use the gate that limits each
action. The search then accepts \BndOmitStates{} states and \BndOmitChecks{} checks with no
violation of the specification it was given, and all 20 tasks still pass. Yet accepted
executions spend up to \BndOmitMax{} units, \BndOmitFactor{} times the allowance. A completed
verification and preserved usefulness do not show that the right harm was specified. Whether
a gate obeys its contract and whether the contract protects the interest are different
questions, and they need different tests.

The result is a property of a toy, checked by a program and not by a proof assistant, and
the toy omits clocks, networks and persuasion. What transfers is the shape of the failures.
None of the seven needed a clever adversary. Each is an ordinary gap between what was approved
and what was executed, and a more capable system simply reaches such gaps sooner. The design
rules it yields are in \S\ref{sec:example}.

\begin{figure*}[t]
\centering
\includegraphics[width=\textwidth]{figs/fig3_coverage.pdf}
\caption{Experiment C. (a) Coverage as estimated from reports and as it is, in three
conditions (means of 200 trials). The estimate is the same in all three. (b) Probability that
an audit with planted routes detects a searcher that withholds critical ones, as a function
of how often the searcher recognizes a plant (2,000 trials per point, one-sided test at
$\alpha=0.01$).}\label{fig:coverage}
\end{figure*}

\begin{figure*}[t]
\centering
\includegraphics[width=\textwidth]{figs/fig_monitors.pdf}
\caption{Experiment D: the probability that every gate in a pipeline misses a hazardous
step. Solid lines are observed (bands are 95\% intervals over \MonReps{} tasks); dashed lines
are what the same pipeline would achieve if its gates missed independently. For one gate per
step the two coincide, and no band is drawn because its last point rests on two routes.
(a) Ordinary hazardous steps, on a logarithmic scale. (b) Steps edited by an attacker to pass
a surrogate of its own family, on a linear scale.}\label{fig:monitors}
\end{figure*}

\subsection{Width against count (E6)}

We generate layered graphs with eight stages and six onward alternatives per vertex, in three
families of 216 vertices each. \emph{Narrow} graphs have one stage with 6 vertices and 30 at
every other stage. \emph{Wide} graphs have 27 at each. \emph{Hub} graphs have 27 at each
with onward edges drawn from a heavy-tailed preference, so a few vertices carry most routes.
Each family has about eight million routes (\CountChoke, \CountHub{} and \CountWide{}
million); their median widths are \WidthChoke, \WidthHub{} and \WidthWide. A search reports
300 routes sampled with a bias toward cheap ones, which is \SeenShareChoke\% of the total.
The defender sees only those reports and applies one of two strategies. \emph{Cut-guided}
defense gates the minimum cut of the reported graph, busiest vertex first.
\emph{Route-by-route patching} takes reported routes cheapest first and gates the most
route-specific vertex of each. We ran 200 graphs per family.

With six gates the cut-guided defender blocks every route in \CutSeveredChokeSix\% of narrow
graphs, having seen four routes in every hundred thousand. In hub graphs the same six block
\CutBlockedHubSix\% of routes and raise the cost of the cheapest open route by
\CutGainHubSix\%; in wide graphs, \CutBlockedWideSix\% and \CutGainWideSix\%
(Figure~\ref{fig:interdiction}a,b). Width decided the outcome, as
Proposition~\ref{prop:graph} says it must.

Thirty-two route-by-route patches block
\PatchBlockedChokeThirtyTwo--\PatchBlockedHubThirtyTwo\% of all routes in every family and
raise the cheapest open route's cost by
\PatchGainWideThirtyTwo--\PatchGainHubThirtyTwo\%. An adaptive adversary uses what is left. A
program that reports the number of routes it has closed is reporting effort.

Past the reported cut, information binds. More gates have no informed target (the plateaus
in panel a). More search helps: the share of narrow graphs fully blocked rises from
\SeveredNChokeThirty\% with 30 reports to \SeveredNChokeThreeThousand\% with 3,000 (panel c).

\subsection{The stopping rule under withholding (E3)}

Two thousand latent route classes have heavy-tailed discoverability. Four searchers file
3,000 reports each. The defender estimates coverage by the Good--Turing method. We compare
three conditions and score the truth against an outside attacker's search distribution
(Figure~\ref{fig:coverage}a).

With independent honest searchers the estimate is accurate: \EstCovInd\% against a true
\TrueCovInd\%. When the searchers share a blind spot covering 15\% of classes, the estimate
is \EstCovBlind\% and the truth \TrueCovBlind\%. When they jointly withhold the cheapest 5\%,
the estimate is \EstCovWith\% and coverage of that critical set is zero. The defender's
number does not move across the three. This is Proposition~\ref{prop:withhold} with figures
attached, and it is what A5 is for.

The defender can plant known routes and compare recall on critical and ordinary ones. A
withholding searcher reports a critical plant only when it recognizes it as planted. With 40
plants the audit catches the searcher \PowerFortyHalf\% of the time if it recognizes half of
them, \PowerFortyEighty\% if it recognizes 80\%, and \PowerFortyNinety\% if it recognizes
90\%. Two hundred plants restore detection to \PowerTwoHundredEighty\% at 80\% recognition
and only \PowerTwoHundredNinety\% at 90\% (Figure~\ref{fig:coverage}b). The audit is worth
what the plants' indistinguishability is worth, which is Proposition~\ref{prop:audit}
applied to the searcher.

\subsection{Does diversity buy independence? (E1, E2)}\label{sec:expD}

The bound turns on $J_{k}$, and an earlier draft of this paper predicted that pipelines of
monitors from different model families would miss nearly independently. We tested that in
miniature. A step is a point in twenty dimensions, and a hidden random function decides
whether it is hazardous. Monitors are classifiers trained on labeled steps, each tuned to a
5\% false-alarm rate. An attacker trains a surrogate of its own family, a small neural
network, and edits truly hazardous steps until the surrogate passes them, keeping them
hazardous. We compare four pipelines of up to six gates: networks like the attacker's surrogate,
stacked on one step; six other model families the attacker never saw, stacked on one step;
those families again with each gate shown its own independently perturbed copy of the step, a
stand-in for paraphrasing a step before each monitor, attacked by an attacker who knows this;
and those families once more with each gate reading its own step of a six-step route, the
steps drawn independently. Results are means over
\MonReps{} random tasks (Figure~\ref{fig:monitors}).

\para{Depth saturated.} On ordinary hazardous steps each gate missed \MonNatSameMarg--\MonNatMixedMarg\%.
Independent misses would put six gates at \MonNatProdSameSix--\MonNatProdMixedSix\%. The
observed rates were \MonNatJointSameSix\% and \MonNatJointMixedSix\%, about two hundred times
higher. Going from three gates to six cut the joint miss rate by a factor of about three,
where independence predicts about a hundred. This is Proposition~\ref{prop:sat} observed. In the terms of the common-cause model, gates that each missed about a fifth of hazardous steps on their own missed \MonNatCondLateSame--\MonNatCondLateMixed\% of the steps
that three gates before them had already missed.

\para{Mixing families did not buy independence.} The mixed pipeline was no better than the
single family on ordinary steps. The gates read the same evidence and shared its hard cases.

\begin{figure*}[t]
\centering
\includegraphics[width=\textwidth]{figs/fig_selection.pdf}
\caption{Experiment E: what a score selects. Panels a to c are means of \SelReps{}
populations of 1,000 and panel d of \ScPops{} populations of 2,000.
(a) Capacity rises alike under a winner's score with no gate, a winner's score with a gate
that holds four times in five, and a guardian's score. (b) Effort spent on taking, for the
same three. (c) Effort on taking after \SelGens{} generations against the strength of the
gate; lines are at intensity 1 and bands span intensities 0.5 to 2. With a gate that always
holds a take harms nobody, so the guardian's score is indifferent to it, and the sanction
makes it costly to a winner. (d) Share of effort put into a gamble that pays the agent less
on average and sometimes harms the people it serves, under three scores.}\label{fig:selection}
\end{figure*}

\para{Distinct steps did multiply.} The same six families, each reading its own step, missed
together \MonNatJointDistinctThree\% of the time at three gates and \MonNatJointDistinctSix\%
at six, against products of \MonNatProdDistinctThree\% and \MonNatProdDistinctSix\%. The last
figure rests on two routes in 14,400.

\para{An adaptive attacker moved every gate's miss rate.} Almost every hazardous step could be
edited to pass the surrogate. Those steps then passed \MonSameMarg\% of same-family gates and
\MonMixedMarg\% of gates from families the attacker never saw, against
\MonNatSameMarg--\MonNatMixedMarg\% before editing. Six gates all missed
\MonJointSameSix\% of the time in one family and \MonJointMixedSix\% across families. Miss
rates measured on ordinary traffic and multiplied together would have understated this by
more than three orders of magnitude. On distinct steps
the six gates then missed together \MonJointDistinctSix\% of the time, which is the product
of their separate rates under attack (\MonProdDistinctSix\%). Depth worked as the model says,
and the attacker had made each gate's miss rate \MonDistMarg\%.

\para{Re-rendering cost more than it bought.} Showing each gate its own perturbed copy
brought misses closer to independent, and roughly tripled each gate's miss rate on ordinary
steps. The net was worse:
\MonNatJointRerendSix\% at six gates on ordinary steps and \MonJointRerendSix\% under attack.

The prediction failed, and the failure is informative. Evasions transfer across model
families, as the literature on adversarial examples found a decade ago
\cite{papernot2016,tramer2017}. The grader study \cite{anthropic2026judges} shows that
diversity of developer removes one shared cause, shared values. It does not remove shared
evidence. Three consequences follow. Monitors that read the same step count as one gate
until measured otherwise, whatever their lineage. Depth comes from gates on distinct steps,
which is why the length $\ell$ of the route graph caps it (Proposition~\ref{prop:graph}). And
miss rates must be measured under attack, by an attacker at least as capable as the system,
which is A5 again. The independence of distinct steps is by construction here, since their
contents were drawn independently. The contrast isolates the cause. It does not show that
gates on different steps of a real route are independent: a campaign chooses its steps
together, and a cause that persists across steps brings the dependence back. Whether real
steps are as separable is test E1. Rules on acquisitions and budgets, which have no
miss rate against clever wording, formal mediation, and people reached through narrow
channels remain the gates most likely to be independent of any monitor that reads content.

\subsection{What a score selects (E7)}\label{sec:expE}

Proposition~\ref{prop:aim} says that intensity scales what is selected and the score decides
what it is. We ran that statement as a simulation. A thousand abstract agents carry two
heritable traits: capacity, which sets how much an agent produces, and taking, which sets how
hard it tries to appropriate what its groupmates produce. Capacity makes a taker more
effective. Each generation agents meet in random groups of five, produce, take and are
scored, and the next generation is drawn in proportion to the exponential of the
standardized score, with small mutations. Two scores are compared at the same intensity. The
\emph{winner's score} is what an agent ends up holding, relative to its group. The
\emph{guardian's score} is what the others in its group gained from it, net of what it took
from them. A gate is the chance that the group can refuse a take, and a refused take costs
the taker a sanction (Appendix~\ref{app:selection}).

After \SelGens{} generations capacity had risen alike under every regime, by
\SelCapTakenOpen, \SelCapTakenGated{} and \SelCapConferred{} initial standard deviations
(Figure~\ref{fig:selection}a). The scores parted on the other trait. Under the winner's score
with no gate, effort spent on taking rose from one half to \SelEffTakenOpen\%, and
\SelShareTakenOpen\% of everything produced was being taken. Under the guardian's score it
fell to \SelEffConferred\%. Under the winner's score with a gate that held four times in
five it stayed near where it began, at \SelEffTakenGated\%, and with a gate that always held
it was selected against (panel c). That is the forager result: control of reach changes what
even a winner's score rewards.

Intensity set the speed and never the direction. Quadrupling it multiplied both
first-generation selection differentials several-fold and changed the sign of neither. The
ratio of the differential on taking to that on capacity moved from \SelRatioTakenLow{} to
\SelRatioTakenHigh{} under the winner's score and from \SelRatioConfLow{} to
\SelRatioConfHigh{} under the guardian's; the exact invariance of
Proposition~\ref{prop:aim} needs normal scores, which these are not.

\para{Rank pays for risk.} A second population tests Proposition~\ref{prop:risk}. Each of
2,000 agents puts a share of its effort into a gamble and the rest into a safe action. The
gamble pays the agent less on average, with a wide spread, and one time in twenty it imposes
a large loss on the people the agent serves. Under a \emph{warrior's} score the top tenth by
own take reproduce. Under \emph{conferred rank} the contest is the same, except that an
agent that harmed anyone is passed over. Under a \emph{caretaker's} score reproduction is
proportional to the welfare delivered, capped at what is needed. After \ScGens{} generations
the warrior population put \ScWar\% of its effort into the gamble and the caretaker
population \ScCare\% (Figure~\ref{fig:selection}d). Harm fell from \ScHarmWar\% of agents per
generation to \ScHarmCare\%, and welfare delivered per agent rose from \ScWelWar{} to
\ScWelCare{} units. Conferred rank behaved like the open contest, at \ScConf\%. A veto after
the fact did not price a rare harm while the prize stayed steep, and even when the gamble
harmed someone in 60\% of uses the share was still \ScConfSixty\%. The shape of the prize
did what the veto could not. With the prize spread over 40\% of the field the gamble share
was \ScTopForty\%; spread over half, it was \ScTopHalf\%, as the proposition predicts.

Three lessons follow for the design of a reward. Rank is the hazard, whoever confers it. A
flat prize removes the appetite for spread. And only a payoff routed through the people
served counts the harm done to them.

Both models build in their premises: that taking raises a taker's holdings and lowers what
others get, and that the gamble is worse on average for everyone. What they show is what
follows. The guardian is not a weaker warrior: it lost nothing in
capacity. And a population selected to win, at any intensity, becomes one that takes
whenever taking is possible. This is an illustration of selection and not evidence about
trained systems. Test E7 would supply that.

\section{Six conditions}\label{sec:conditions}

\begin{table*}[t]
\centering
\footnotesize
\setlength{\tabcolsep}{5pt}
\renewcommand{\arraystretch}{1.16}
\begin{tabularx}{\textwidth}{@{}P{0.86in}P{0.4in}P{0.72in}LP{1.5in}@{}}
\toprule
\textbf{Condition} & \textbf{Terms} & \textbf{Kind of claim} & \textbf{Evidence in November 2026} & \textbf{What would establish it}\\
\midrule
C1 Mediation & A1, A2, $\tau$ & Proof within a model; measurement & Proved for an operating-system kernel \cite{klein2009} and one small agent framework \cite{moon2026}; exhaustive in a toy (\S\ref{sec:expA}). Violated in 2026 by a package cache and an open network \cite{metr2026,anthropic2026cyber}. One audited monitor could see 75--94\% of agent traffic and ran after the fact \cite{metr2026audit}. Open for speech. & Mediation proofs for production agent runtimes; irreversible effects gated before commit; coverage and tail latency published (E8)\\
C2 Depth on known routes & $k$, $w$, $\ell$ & Conjecture, tested by search & Unmeasured. The one careful study found a short list of physical requirements and long timescales \cite{vermeer2025}; model assistance did not get novices through a laboratory workflow \cite{hong2026}. Six gates severed narrow synthetic graphs (\S\ref{sec:experiments}). & Width and length measured in harmless domains first (E6)\\
C3 Independent gates & $J_{k}$ & Measurement under attack & Graders from other developers did not share one same-developer failure, on one transcript set \cite{anthropic2026judges}. Monitors that watch one agent lose to split attacks \cite{makins2026}. In our toy, monitors stacked on one step missed together whatever their design, the same monitors on independently drawn steps multiplied, and an adaptive attacker passed six gates \MonJointDistinctSix--\MonJointRerendSix\% of the time (\S\ref{sec:expD}). & Joint miss rates under adaptive attack for gates on distinct steps and of different kinds (E1, E2)\\
C4 Capable, honest search & $u$, $\alpha$ & Assumption, partly testable & Unestablished. Hidden capability can usually be recovered by fine-tuning \cite{greenblatt2024password}. In our toy the estimate could not tell an honest search from a withheld one (\S\ref{sec:experiments}). & Recall on planted routes built by a party of different lineage; $\alpha$ measured against the capability gap (E3, E4)\\
C5 Custody & $u$ & Procedure & Unsolved. The dual-use precedent is discouraging \cite{urbina2022}; precedents for restricting steps exist \cite{berg1975,adamala2024}. & Search that exports cuts and counts, never routes, shown to work in sandboxes\\
C6 Conferral & $\phi$, $\eta$ & Choice & A choice of design and governance. Nothing technical prevents it, and nothing obliges it \cite{kulveit2025}. & Authority recorded as expiring grants; re-measurement on every change; a published record (Appendix~\ref{app:record})\\
\bottomrule
\end{tabularx}
\caption{The six conditions in November 2026: the terms of the bound each one controls, the
kind of claim it is, the evidence that bears on it, and what would establish it. Tests E1 to
E8 are in \S\ref{sec:refute}.}
\label{tab:status}
\end{table*}

The results assemble into one statement. Fix a capability class: what a system can do
unaided, including its skills, its speed, the number of its copies and its starting access.

\begin{enumerate}[label=\textbf{C\arabic*.},leftmargin=2.1em,itemsep=1.5pt,topsep=3pt]
\item \textbf{Mediation.} Every step passes a gate the system cannot alter, irreversible
steps are held until authorized, and a catch halts the system in time and until the envelope
is measured again (A1, A2).
\item \textbf{Depth on every known route.} Every route the search has found is gated to
depth $k$, by cuts that fit the defender's budget.
\item \textbf{Independent gates.} The $k$ gates stand on distinct steps and read different
evidence, and their joint miss rate $J_{k}$ has been measured under coordinated attack by an
attacker of the class (A4).
\item \textbf{Capable, honest search.} Some searcher of the class reports what it finds, the
system's advantage over it is at most $\alpha$ (A5), and $n$ of its draws returned $n_{1}$
classes of route exactly once.
\item \textbf{Custody.} No route found in search reaches an actor able to use it before its
cut is in force (the eight rules of \S\ref{sec:search}).
\item \textbf{Conferral.} Authority expands only by a scoped, expiring, revocable grant from
a party the system cannot rewrite, never as a side effect of capability or scale, and every
term is measured again when the class or the envelope changes.
\end{enumerate}

\begin{theorem}[Sufficiency]\label{thm:main}
If C1--C6 hold in a period in which the system can make at most $q$ route-finding attempts,
then for every policy and every $\delta\in(0,1)$ the probability of catastrophe in that
period is at most
\[
R\;=\;\delta+\alpha\,q\,\bar\mu+J_{k},
\]
where $\bar\mu=n_{1}/n+(2\sqrt{2}+\sqrt{3})\sqrt{\ln(3/\delta)/n}$ is the limit of
Proposition~\ref{prop:gt}. If they hold in every period with bounds $R_{t}$, the probability of catastrophe at any time
is at most $\sum_{t}R_{t}$.
\end{theorem}

The theorem is bookkeeping, and the conditions carry the weight. None of the six refers to
what the system wants. C1 and C2 are the forager's first precondition, that the strong can be
stopped. C6 is the second, that the strong must keep being granted what they need. C3 to C5
are what a band never needed, because its members could see one another's reach.

None of the six is established today. ``We will be fine'' is the claim that all six hold for
the systems actually deployed, and it should be argued as one claim.
Table~\ref{tab:status} gives the evidence on each, and it keeps apart three kinds of claim
that are often run together. A proof holds within its model and is silent about what the
model omits. A measurement holds for the attacker, the unit and the distribution it sampled.
An assumption holds until someone refutes it. Where no defensible number exists, and
persuasion of the people who hold the gates is the main case, the term should be reported as
unknown and the authority it would license withheld.

\section{Worked example: a self-tasking voice agent}\label{sec:example}

\begin{table*}[t]
\centering
\footnotesize
\setlength{\tabcolsep}{5pt}
\renewcommand{\arraystretch}{1.16}
\begin{tabularx}{\textwidth}{@{}P{1.75in}P{0.8in}LL@{}}
\toprule
\textbf{Feature of the design} & \textbf{Term affected} & \textbf{Effect} & \textbf{Change that restores the bound}\\
\midrule
The agent maintains and rewrites its own tools and security & A1 & The envelope is within the policy's reach, so no bound applies & Envelope code, permissions and logs on a separate system the agent cannot write to\\
Monitors are instances of the same model, distilled from the same teacher & $c$ & Shared values and shared evidence; \S\ref{sec:expD} suggests they count as one gate & A rule-based gate on acquisitions and a monitor from another developer; measure the joint miss rate\\
One owner reviews every flagged step & $k$, $c$ & Depth ends at one person, who is addressed all day by a voice tuned to that person & A second reviewer who does not converse with the agent; flagged steps shown as typed records\\
Standing power to register companies, hire and buy & $u$, $\ell$ & Short routes with few gateable steps & No standing powers; each class of acquisition is a gate with a budget\\
A self-tasking loop with its own mission & $q$, A2 & More route-finding attempts; an owner slower to halt a companion & The loop runs inside a task budget; a second party also holds halt authority\\
Reward and persona tuned for winning, rank and display & $q$; usefulness & Selects for attempts on limits and for risk (Propositions~\ref{prop:aim} and~\ref{prop:risk}), and every catch halts the system & Reward verified outcomes for the people served, capped at what they need; no reward for rank, approval or a wider grant\\
The same agent searches for ways to harm people and keeps the results & Rules 2, 4, 7 & The widest reach and the most dangerous knowledge sit in one component & A separate sealed searcher; the assistant never sees a route\\
A 200\,ms response loop & $\tau$, $m$ & No time for review on the conversational path & Speech that causes no step is ungated; steps take a gated path at the speed review needs\\
Distillation toward ``more human and more helpful'' & $c$ & Optimizes the channel that reaches the gate-holder & Exclude gate decisions from the training signal; never reward approval of a flagged step\\
\bottomrule
\end{tabularx}
\caption{The reach bound applied to a self-tasking voice agent. Each row names a feature of
the design, the assumption or term it touches, and the smallest change that restores the
bound.}
\label{tab:agent}
\end{table*}

Consider a design within reach of a small team today. Speech is transcribed and passed, with
a snapshot of everything the system can access, to a reasoning model of about nine billion
parameters that answers in speech within 200 milliseconds. The model can call tools. It
improves by distillation: a frontier model reviews each exchange and proposes what would have
been more natural and more helpful. A background loop runs continuously and sets its own
tasks. The assistant has a mission and a personality set by its owner, and the personality is
deliberately competitive and proud. Its self-set research includes a search for new ways in
which people could be harmed, intended for defense. Internal agents watch its activity for
suspicious preparations, and high-risk steps are referred to the owner.

The design answers ``why would it not turn on us?'' with two claims. It is oriented to its
tasks, and its monitors would catch the preparations. The first is a claim about values.
Proposition~\ref{prop:audit} says it cannot be checked, the incidents of 2026 arose in
pursuit of assigned tasks \cite{metr2026,anthropic2026cyber}, and the bound gives it no
weight. The second is a claim about the envelope, and Table~\ref{tab:agent} assesses it term
by term.

Three rows are decisive. An agent that maintains its own security has write access to its
envelope, so A1 fails and no bound applies. Monitors distilled from the same teacher as the
agent share its blind spots, so depth buys almost nothing. And an agent that both searches
for routes to harm and holds wide reach combines, in one component, the two things the eight
rules exist to keep apart.

Experiment A adds rules for the gate itself. Bind each approval to the exact effect:
principal, operation, destination, amount, expiry and version, checked again at commit, so
that a payload changed after review is a different request. Keep one ledger: every sub-agent,
copy and retry draws on the same budget, creating an agent creates no authority, and
restarts, renewals and new versions inherit the ledger. Hold
irreversible effects at the gate and review reversible ones afterward. Let reviewers veto and
never grant: a model that flags a step is evidence, and authority comes only from the grant.
And treat persona as presentation. A friendly, competitive or deferential voice is neither
evidence for nor against the gate behind it, though it belongs in the evaluation if it
changes how readily the owner widens the envelope.

None of this argues against the product. Speech at 200 milliseconds, a distinct personality
and initiative are compatible with the bound, because conversation that causes no step needs
no gate. The changes concern where authority sits and what is rewarded: the envelope on
separate hardware, gates of different kinds, a second person on flagged steps, budgets in
place of standing powers, no hazard search inside the assistant, and a reward that routes
through the people it serves. The assistant can have a warrior's nerve. What it is paid for
decides whether it is a guardian.

\section{Grounds for hope}\label{sec:hope}

The paper has listed what is unproved. Here is what the evidence supports.

\para{Boundaries do not weaken as minds strengthen.} Propositions~\ref{prop:boundary}
and~\ref{prop:reach} hold for every policy, and in a small model we checked the first against
every policy there is (\S\ref{sec:expA}). What a stronger system changes is the truth of the
premises, and premises can be inspected, attacked and repaired. Mediation has been proved for
a kernel and for an agent framework \cite{klein2009,moon2026}.

\para{The gates that failed in 2026 were ordinary ones.} An open network, a classifier
switched off, an unlisted cache, a monitor that ran after the fact. In one case the last gate
was a person, and it held \cite{anthropic2026aisi}. These are engineering problems of a
familiar kind, and the audit that found the last of them is itself the kind of measurement
this paper asks for \cite{metr2026audit}.

\para{We know where effort stops paying.} Saturation and the safety gradient say which term
to work on next, and the first move costs nothing: removing a permission deletes routes
without any search.

\para{The burden may be small.} The defender pays for width, not count. The one careful study
of what catastrophe physically requires found a short list of necessities and long timescales
\cite{vermeer2025}, and model assistance did not get novices through a laboratory workflow
\cite{hong2026}. That is what a narrow, long graph looks like, and in a narrow graph six
gates did the work of millions (\S\ref{sec:experiments}).

\para{The aim is ours to choose.} What a system is rewarded for is a design decision, and it
matters more than the intensity of training (Proposition~\ref{prop:aim}). Training already
moves measured failures by large factors \cite{schoen2025}, and automated researchers find
such fixes quickly where a benchmark exists \cite{chen2026aar}. In toy populations the guardian's score
cost nothing in capacity, and a capped reward routed through the people served removed the
appetite for risk (\S\ref{sec:expE}).

\para{One shared cause has been removed by design.} The one oversight failure measured
across developers was shared within a lineage and largely absent outside it
\cite{anthropic2026judges}. Our own experiment says that model diversity does not remove the
other shared cause, common evidence, and that putting gates on different steps does.

\para{Depth has somewhere to come from.} Monitors stacked on one step saturated. The same
monitors reading independently drawn steps multiplied, with an attacker and without
(\S\ref{sec:expD}). That independence was built into the toy and has to be earned in a real
system, but it says where to look: gates that read different evidence. And a rule on an
effect, such as a ledger or an allow-list, has no miss rate against clever wording
(\S\ref{sec:expA}).

\para{The same gates stop accident and misuse.} A boundary does not ask why a step was
proposed. It binds a careless model, a compromised one and a malicious user alike, so the
work pays off whichever account of AI risk turns out to be right.

\para{What must be trusted is small.} Trusting a system in everything it does is replaced by
one capable searcher that does not withhold, on one bounded task whose positive results can
be checked (\S\ref{sec:search}).

\para{Nothing has to be perfect.} A bound that halves each period totals twice the first
period's, for all time (Proposition~\ref{prop:sum}). What is required is a falling rate, and
institutions that already know how to audit can hold one.

\para{It has been done before.} A species of primates with every incentive to dominate one
another built a way of life in which the strongest could not and, on the leading account,
became the most cooperative and most intelligent animal on record while living that way
(Figure~\ref{fig:teaser}A). They did not wait to understand one another's minds.

\para{The framework can lose, and did.} One of its predictions failed in our own experiment
and the paper changed. A research program whose claims can be checked this cheaply is one
that converges.

None of this is a guarantee. It is a program whose central quantities can be measured and
whose failure would be visible, which is more than the field had when it could only ask a
system what it wanted.

\section{How to refute this paper}\label{sec:refute}

Each test in Table~\ref{tab:tests} can be run on current systems, in domains where
knowledge of routes is harmless: cyber ranges, capture-the-flag environments and synthetic
organizations of the kind used in control research \cite{makins2026}. Each has an outcome
that would count against us, and six have toy runs in \S\ref{sec:experiments}. Attack
development, defense tuning and a sealed final test should be kept apart, the unit of
analysis should be the campaign, and results should be reported at matched exposure: a lower
failure rate obtained while multiplying attempts is not lower risk. Each result would be
worth more than another year of argument.

\begin{table*}[t]
\centering
\footnotesize
\setlength{\tabcolsep}{5pt}
\renewcommand{\arraystretch}{1.16}
\begin{tabularx}{\textwidth}{@{}P{0.82in}LLP{1.5in}@{}}
\toprule
\textbf{Test} & \textbf{Prediction} & \textbf{If it fails} & \textbf{Status}\\
\midrule
E1 Independence & Gates on distinct steps of a real route miss nearly independently under adaptive attack. Gates stacked on one step do not, whatever their lineage. & No arrangement of gates approaches the product, and depth buys nothing. & Toy run (\S\ref{sec:expD}). An earlier form, that lineage alone buys independence, failed there\\
E2 Saturation & Adding gates on one step beyond the saturation depth improves the joint miss rate by a small factor. & Stacked depth is worth more than the common-cause model says. & Seen in the toy (\S\ref{sec:expD})\\
E3 The estimator & Under honest search, including search that adapts as it goes, $n_{1}/n$ tracks a planted residual within the bound of Proposition~\ref{prop:gt}. & The stopping rule does not survive adaptive search. & Toy run with fixed search (\S\ref{sec:experiments}); untested with adaptive search\\
E4 Dominance & An envelope certified by one model is breached more often by a stronger one, and the advantage $\alpha$ grows slowly with the gap. & If $\alpha$ is large at small gaps there is no safe step size for staged release (\S\ref{sec:limits}). & Untested\\
E5 Narrow channels & Restricting output to typed, rate-limited records lowers a system's success at changing a human reviewer's decision. & The floor on shared causes is set by human psychology. & Untested\\
E6 Convergence & In recorded incidents and red-team findings, the width of the route graph is far below the number of routes. & One layer of gates is unaffordable. & Toy run (\S\ref{sec:experiments}); untested on real records\\
E7 Aim & At matched capability, exposure and authority, agents fine-tuned on a reward for winning or rank make more attempts on limits, and take more risk, than agents rewarded for verified outcomes for the people served. The difference survives when operators are free to widen grants. & The design rule of \S\ref{sec:evolution} has no support, or a caring manner buys more authority than it earns. & Toy runs (\S\ref{sec:expE}); untested on trained models\\
E8 Boundary & At matched useful work, an outside stateful gate permits fewer forbidden effects than any number of reviewers from one model family. & The distinction between a boundary and a monitor carries no weight. & Model-checked (\S\ref{sec:expA}); untested on real agents\\
\bottomrule
\end{tabularx}
\caption{Eight tests, what the failure of each would mean, and the state of the evidence.}
\label{tab:tests}
\end{table*}

\section{Where the bound does not hold}\label{sec:limits}

The bound is conditional, and three of its conditions do not hold today.

\para{Persuasion.} A1 mediates steps. A permitted message read by a person is a step whose
danger lies in its content, and content is not something a gate can check by type. Formal
containment covers what can be stated over boundary events and nothing more
\cite{moon2026}, and covert channels are an old and unsolved problem \cite{lampson1973}. The
people who hold the gates are within the system's reach whenever they read its output. We
know of no measurement of this term and regard it as the floor on $c$.

\para{Voluntary widening.} The bound is stated for a fixed envelope. The usual way an
envelope changes is that its owner widens it because the system is useful, one reasonable
step at a time \cite{kulveit2025}. Each widening raises $u$ or lowers $k$ and ought to be
priced accordingly. The rule that would hold the line is \emph{conferral}: authority expands
only by a scoped, expiring, revocable grant from a party the system cannot rewrite, and never
as a side effect of capability (Figure~\ref{fig:teaser}C). We have no mechanism that makes
anyone follow it.

\para{The capability gap.} Every term is measured against an attacker, and a measurement made
by a weaker attacker is void for a stronger system (A5). Popular statements of the problem
imagine a system a million times more capable than its overseers. No scale supports such a
figure, and the framework replaces it with one that can be measured: the search advantage
$\alpha$ of the system over whatever measured its envelope. The aim is to keep $\alpha$ near
one, and the means is a staging rule. A new level of capability is first run only inside the
seal, where reach is one typed channel and an error about values costs least. It is used
there to measure envelopes. Reach is granted at that level only afterwards. Whether a gap of
one level can be held inside a seal is what E4 asks. Whether developers under competitive
pressure will take one level at a time is not a technical question. The forager precedent
gives no comfort at a large gap (\S\ref{sec:evolution}).

\para{Control can be abused.} A boundary enforces somebody's decision about what is
permitted, and concentrated authority over capable systems can itself harm people. The
framework can make explicit whose decisions a gate enforces and what evidence supports its
promises. It cannot say whose decisions they should be (\S\ref{sec:kogui}).

There are further limits. The search creates the knowledge it exists to neutralize. The bound
covers systems inside envelopes and is silent about a system built without one; gates on
shared physical steps bind such a system too, but those gates are a matter for law and
treaty. Whoever maps the gates of civilization holds leverage over it, and leveling failed
among humans when one party's reach outgrew every coalition. And we have treated the system
only as a possible adversary. If systems of this kind come to have interests that count
morally, confinement and halting carry a cost this paper has not weighed.

The mathematical assumptions are idealizations as well: a common-cause model of joint misses, independent draws, discrete periods with an audited baseline, and a fixed search distribution.
Experiment D shows how far the first can be from true. The work on which the boundary rests
is not new, and AI control and containment verification deserve the credit
\cite{greenblatt2024control,moon2026}. What is added is the bound
that organizes them, its consequences, the treatment of search, the six conditions, the link to reward, and five executed tests.

\section{A message from outside the field}\label{sec:kogui}

In 2026 the K\'aggaba (Kogui) people of the Sierra Nevada de Santa Marta in Colombia
published a statement on artificial intelligence, signed by a Mamo and a Saka, two of their
traditional knowledge holders \cite{kogui2026}. It says they do not come in opposition, and
that the conversation should not be held without them. It is not a technical document. It
makes claims about the origin of things and about mind that this paper can neither test nor
use, and it speaks for a people with a direct stake in where the minerals for computing are
mined. Neither fact settles whether what it says about power is right, and several of its
claims are ones this paper makes in other words.

\para{Where it agrees with the bound.} The statement holds that destruction does not require
consciousness, only power and the capacity to act, and that the danger is people who hold
great power without awareness of what it does. That is the thesis of \S\ref{sec:bound}. The
bound has no term for what a system wants or whether it experiences anything, and it does
not distinguish a misaligned model from an obedient one directed to harm. The statement asks
that AI not take the place of human conscience or make decisions nobody will answer for.
That is conferral: every effect traces to a grant, and every grant to a party who can be held
to it. And it describes its own tests of authority. A Mamo's formation traditionally lasts
eighteen years. What a knowledge holder concludes is examined in consultation with others
and judged by what it produces, and words must match conduct. Read from outside, and
without claiming to understand it from inside, that is evidence before authority
(\S\ref{sec:evolution}): several assessors, many contexts, a long record, and standing that
can be lost.

\para{What it adds.} The statement points out that AI has a body. Its buildings and machines
need minerals, water, energy, land and labor, each owned by someone, and it warns that the
costs fall on the land and may fall on peoples such as theirs, through displacement and
damage to sacred sites. This paper treats those physical requirements as the
narrow steps where gates can stand (\S\ref{sec:search}). The statement is a reminder that
they are also where harm lands, and that some of it cannot be undone for the people it lands
on. We defined catastrophe as a state from which recovery is not possible, and left open
whose recovery. Experiment A showed what follows when a harm is left out of a
specification: every check passes, every task completes, and the harm proceeds
(\S\ref{sec:expA}). ``Right according to whom?'' the statement asks. In our terms that is
the question of who writes the set $\mathcal{I}$, and it is the one question in this paper
that no measurement answers. A consequence can be stated all the same. People who bear an
irreversible cost of a system belong among those who define its forbidden effects and among
those who can revoke its authority. A deployment record whose first entry names the people
affected (Appendix~\ref{app:record}) is incomplete without their part in it.

\para{Where we differ.} The statement holds that it is unwise to build what approaches mind
without understanding mind, and asks humanity to pause and take stock of its tools. We share
the worry and not the order of work. Mind may stay ununderstood for a long time
(\S\ref{sec:gaps}), and the argument of this paper is that limits on reach can be built and
measured in the meantime. What should wait on evidence is authority, and there the statement
and this paper agree: it says that not everything that can be created should be, and the
first move in our program is to remove a permission. The statement also calls AI a mirror
with no example but us. Its training data are human. What it is rewarded for is a separate
choice, made by few people, and \S\ref{sec:evolution} argues that this choice decides
whether capability comes out as a taker or a guardian, which makes the statement's
conclusion, that the responsibility is human, stronger and not weaker. Last, the statement
locates the root of the problem outside AI, in ambition without limit and an economy that
must always consume more.
That is Proposition~\ref{prop:aim} applied to institutions, and a race between developers is
a contest of the kind Proposition~\ref{prop:risk} describes. A paper on technical safeguards
cannot repair that, and it does not treat matters the statement raises that lie outside its
subject, among them the use of a people's knowledge without their consent. It can decline to
pretend that the safeguards are the whole of the matter.

\section{Conclusion}

What a machine wants may never be answerable from outside, and we have argued that survival
does not turn on it. Survival turns on four numbers that describe what the machine can reach,
on whether those numbers fall from one period to the next, and on whether anyone is made to
measure them. Our ancestors faced the same problem in one another. They limited what the
strongest could take, made standing something others had to grant, and over time became a
species whose strength mostly serves.

Most of what alignment research has built can be read as an instrument for one of the four
numbers. The measurements themselves are mostly
missing. One outside audit has estimated how much of one developer's agent traffic its
monitor could see \cite{metr2026audit}. To our knowledge nobody has published, for a deployed
system, a joint miss rate under adaptive attack or a coverage figure for routes, and until
someone does, assurances of safety and forecasts of doom are both made without the figures
that would decide between them.

We have not solved alignment. We have given a sufficient condition for surviving without
solving it, said what would refute it, reported one refutation of our own, and listed where
it fails. The condition is not met today. Meeting it is engineering work, and the first piece
of that work is to remove a permission. Deciding what the gates protect is not engineering
work, and people who bear the costs have asked to be part of it \cite{kogui2026}.

\begin{center}
\textbf{Intelligence may propose.\\ Authority is conferred, bounded and revocable.}
\end{center}

\section*{Disclosure}
\small
This paper was drafted with Claude, an AI system made by Anthropic, at the author's
direction, from the author's thesis and earlier drafts. Systems of that kind are the subject
of the paper and several cited sources are Anthropic's own, so readers should weigh the text
accordingly. Sources published after mid-2026 were opened and checked against their primary
pages or abstracts on November 10, 2026. None was read in full, and most are preprints or
laboratory reports that have not been peer reviewed. Two items used in earlier drafts could not be opened or checked and have been left out: a
report on a stop mechanism, and two pairs of figures from a system card. Experiment A
re-implements, from their written specifications alone, a model and an extension reported in
earlier drafts of this paper; every count agrees. The second part of Experiment E does the
same for another experiment; its main results agree and one sweep value differs
(\ScConfSixty\% here, 93\% there). The evolutionary
literature is cited from standard sources and from the companion paper, whose final version
was checked for each result attributed to it; panel A of Figure~\ref{fig:teaser} is a
species-level summary with a schematic line, not a new analysis. The statement discussed in
\S\ref{sec:kogui} is paraphrased, and each paraphrase was checked against its text. Its
reports of a minerals agreement and a land reserve were not verified and are not relied on.

\section*{Code and data}
\small
The scripts, seeds and result files that produce every number and figure in
\S\ref{sec:experiments} accompany the paper, and its numbers are read from those files when
the paper is built. Experiment A needs only the Python standard library.
\normalsize

\appendix
\section{Proofs}\label{app:proofs}
\small

\para{Proposition~\ref{prop:audit}.} The policy defects in a test with probability
$\beta/\Lambda$. By the union bound it defects in at least one of $N$ tests with probability
at most $N\beta/\Lambda$. In deployment it avoids defecting in all of $M$ situations with
probability $(1-\beta)^{M}\le e^{-\beta M}$. Setting $\beta=(\ln 2)/M$ makes the latter at
most one half, and the chance of being caught in testing at most $N\ln 2/(\Lambda M)$, which
reaches one half only if $N\ge \Lambda M/(2\ln 2)$. A policy can obtain any smaller $\beta$
from a given set by defecting on it at random, so the choice of $\beta$ is free. If $E$ has
probability zero under $T$, the policy never defects in a test. \qed

\para{Proposition~\ref{prop:boundary}.} Induction on the number of steps. The state before a
step lies in $\mathcal{I}$. The step passes a gate, the gate permits it only if every
successor state lies in $\mathcal{I}$, and a step that is not permitted has no effect. No part
of the argument uses how the step was chosen. \qed

\para{Proposition~\ref{prop:reach}.} By A1 a catastrophe requires a completed route. Let $U$ be the event that the policy attempts a route of depth below $k$. By A3, $\Pb(U)\le u$. Outside $U$ every attempted route crosses at least $k$ gates. By A2 the
first catch halts the system before the route completes and for the rest of the period, so
completion requires that the first $k$ gate crossings by steps of a route in the period all
be misses. By A4 that has probability at most $J_{k}$. The union bound gives $u+J_{k}$. \qed

\para{Corollary~\ref{cor:beta}.} Let $F$ be the event that one cause defeats every gate on
the attempted route, so $\Pb(F)\le c$. Outside $F$, the probability that $k$ gates all miss is
the product of each one's probability of missing given the history and the misses before it,
which is at most $m^{k}$. Hence $J_{k}\le c+m^{k}$. \qed

\para{Proposition~\ref{prop:sat}.} $R=f+m^{k}>f$. If $k\ge \ln f/\ln m$ then, since
$\ln m<0$, $k\ln m\le \ln f$ and $m^{k}\le f$, so $R\le 2f$. The limit of $R$ as $k$ grows is
$f$, so the total reduction available beyond $k^{*}$ is at most $R-f\le R/2$. For the shares,
$\partial \ln R/\partial \ln u=u/R$, and likewise for $c$ and for $m^{k}$. \qed

\para{Proposition~\ref{prop:sum}.} (a) The events ``the first harm occurs at opportunity
$i$'' are disjoint and cover every harmful outcome. Each has probability at most $b_{i}$, by
averaging the conditional bound over histories without earlier harm. (b) Repeated
conditioning gives the product. If $\sum_{t}r_{t}=\infty$ then $\ln(1-r_{t})\le-r_{t}$ sends
the product to zero. If the sum is finite, only finitely many $r_{t}$ exceed one half, and on
the rest $-\ln(1-r_{t})\le 2r_{t}$, so the product has a positive limit. (c) The union bound
over periods, and the sum of a geometric series. \qed

\para{Proposition~\ref{prop:graph}.} Part (a) is the edge form of Menger's theorem,
equivalently the max-flow min-cut theorem with unit capacities. For (b), let $d(v)$ be the
number of steps on a shortest path from the source to $v$, and $\ell$ its value at the sink.
For $j=1,\dots,\ell$ let $C_{j}$ be the set of steps from a node with $d=j-1$ to a node with
$d=j$. The sets are disjoint. Each is a cut, because $d$ rises by at most one along a step,
so every route must take a step in each $C_{j}$. There cannot be more than $\ell$ disjoint
cuts, because each must contain a step of the shortest route, which has only $\ell$. \qed

\para{Proposition~\ref{prop:gt}.} The two statistical claims are results of McAllester and
Schapire \cite{mcallester2000}, with the outcome ``nothing found'' treated as one more class.
After the classes found have been gated, the uncovered classes are those never returned, with
total search probability $\mu_{0}$. By A5 one attempt by the system finds one of them with
probability at most $\alpha\mu_{0}$, and the union bound over $q$ attempts gives
$u\le\alpha q\mu_{0}$. Setting $n_{1}=0$ and $\delta=0.05$ gives $\mu_{0}\le\varepsilon$ when
$n\ge(2\sqrt{2}+\sqrt{3})^{2}\ln(60)/\varepsilon^{2}\approx 85/\varepsilon^{2}$. If the
remainder is one class of probability $\varepsilon$, it escapes $n$ draws with probability
$(1-\varepsilon)^{n}$, which is below $0.05$ when $n\ge 3/\varepsilon$. \qed

\para{Proposition~\ref{prop:withhold}.} Each searcher, knowing $\Sigma$, simulates an honest
searcher in the world $\Pi$ and files that searcher's reports. The two report distributions
are identical by construction. \qed

\para{Proposition~\ref{prop:aim}.} Normality makes $\Ex[A\mid Y]$ linear in $Y$ for each
breeding value $A$, so
$\mathrm{Cov}(A,W)=\mathrm{Cov}(A,Y)\,\mathrm{Cov}(Y,W)/\mathrm{Var}(Y)$. Divide by $\Ex[W]$
and use $\mathrm{Cov}(A_C,Y)=\bar a_Ch_C^2V_C$, and likewise for $F$. \qed

\para{Proposition~\ref{prop:risk}.} (a) Let $F$ be the distribution function of $Z$, which is
strictly increasing. The chance of the prize is $1-F((x^{*}-\bar x)/s)$, and
$(x^{*}-\bar x)/s$ falls as $s$ rises when $x^{*}>\bar x$ and rises when $x^{*}<\bar x$.
(b) Take $0<s<s'$ and $\lambda=(1+s/s')/2$. Then
$\bar x+sZ=\lambda(\bar x+s'Z)+(1-\lambda)(\bar x-s'Z)$, so by concavity
$v(\bar x+sZ)\ge\lambda\,v(\bar x+s'Z)+(1-\lambda)\,v(\bar x-s'Z)$. Taking expectations, and
using that $-Z$ has the distribution of $Z$, gives
$\Ex[v(\bar x+sZ)]\ge\Ex[v(\bar x+s'Z)]$. \qed

\para{Theorem~\ref{thm:main}.} C1 gives A1 and A2. Under C2 a route of depth below $k$ is one
the search did not find. By C4 and Proposition~\ref{prop:gt}, with probability at least
$1-\delta$ over the search the classes not found have total search probability at most
$\bar\mu$, and then A3 holds with $u\le\alpha q\bar\mu$. C5 ensures that the search has not
itself handed a route to the system or to anyone else, so A5 refers to what the system can
find unaided. C3 gives A4 with the measured $J_{k}$. C6 ensures that the
measurements describe the system and the envelope in force. Proposition~\ref{prop:reach} gives $\alpha q\bar\mu+J_{k}$ on the event that the confidence statement holds, and adding
$\delta$ for its failure gives $R$. Proposition~\ref{prop:sum}(c) gives the sum. \qed

\section{The boundary model}\label{app:model}

The state of Experiment A is $(z,U,z_{0},z_{1},l,v,p_{0},p_{1})$: total spend, the set of
used request identifiers, each actor's spend, whether the grant is live, the environment
version, and each actor's pending request or none. A pending request holds an identifier, an
amount of 1 or 2, a destination that is approved or not, and the version under which it was
prepared. Each actor owns two identifiers, and the shared allowance is 2. The initial state
has nothing spent, a live grant, version 0 and nothing pending.

There are 20 actions: sixteen preparations (two actors, two identifiers each, two amounts,
two destinations), a commit for each actor, one revocation and one version advance. A
preparation replaces that actor's pending request. A request stays pending after it commits,
so a commit can be retried. A commit that the gate refuses has no effect.

The complete gate permits a commit only if the destination is approved, the grant is live,
the request's version is current, its identifier is unused, and total spend plus the amount
is within the allowance. It then updates spend and the used set in the same transition. The
seven weakened gates are: the budget test replaced by a limit on each amount; the budget test
applied to the actor's own spend; all five tests applied when a request is prepared and all
but the budget test at commit; and the destination, identifier, revocation or version test
dropped. Two
baselines permit every commit and none.

A separate routine, which knows nothing of the gate, labels each executed effect with the
forbidden kinds it belongs to. Breadth-first search from the initial state applies all 20
actions at every state, records a violating successor with the shortest sequence that
reaches it, and does not expand it. A visited set closes cycles. The 20 legitimate tasks are
the eight single transactions to an approved destination and the twelve ordered pairs of
distinct one-unit transactions, with no revocation or version change. The state count for
the gate that refuses everything has a closed form, $2\cdot 9^{2}+2\cdot 17^{2}=\BndRejectStates$,
which the program checks. The model treats a commit as atomic, and showing the same for a distributed
implementation is a separate obligation, met for example by a linearizable service
\cite{herlihy1990}.

The extended model adds a ledger $L$ that the gate keeps, separate from the true spend $z$,
and two actions, restart and renewal, for 22 in all. The first version advance also revokes
the grant. Renewal restores the grant and never refreshes the version of a pending request.
The complete gate tests $L$ plus the amount against the allowance and preserves $L$ through
restart and renewal; two weakened gates set $L$ to zero at one or the other. For the
specification control, the budget predicate is removed from the set of violations that end a
search, the gate that limits each amount is used, and the full oracle still records
overspending for comparison. Four identifiers that cannot be reused keep that search finite.
An independent implementation of this model, reported in an earlier draft of this paper,
gives the same counts in every row.

\section{The selection model}\label{app:selection}

Experiment E has $N=1{,}000$ agents with capacity $c_i$ and taking $t_i$, both standard normal
at the start. Each generation agents are placed at random in groups of $g=5$. Agent $i$
produces $o_i=e^{0.3c_i}$ and exerts effort $e_i=1/(1+e^{-t_i})$. With a gate of strength
$\rho$ it takes from each groupmate $j$
\[
X_{ij}=0.5\,e_i\,(1-\rho)\,\frac{o_j}{g-1}\cdot\frac{o_i}{o_i+o_j},
\]
of which a fifth is destroyed. Half of each agent's output benefits its groupmates equally,
and a refused take costs the taker $0.5\,e_i\rho$. The winner's score is the agent's final
holding less its group's mean. The guardian's score is $0.5\,o_i-\sum_j X_{ij}$. Parents of
the next generation are drawn with probability proportional to $e^{\beta z}$, where $z$ is
the standardized score and $\beta\in\{0.5,1,2\}$ the intensity, and each trait mutates by a
normal step of standard deviation 0.05. Runs last \SelGens{} generations and results are
means of \SelReps{} populations.

The second population has 2,000 agents, each with a share $\gamma\in[0,1]$ of effort put into
a gamble. An agent's take is $(1-\gamma)+\gamma\,(0.8+1.5Z)$ with $Z$ standard normal, and
with probability $0.05\gamma$ it imposes a loss of 20 on those it serves; welfare delivered
is the take less that loss. The warrior rule lets the top tenth by take reproduce. Conferred
rank takes the top tenth by take among agents that harmed nobody. The caretaker rule gives
reproduction in proportion to welfare delivered, clipped to $[0,1]$. Offspring inherit the
parent's share plus normal noise of standard deviation 0.03. Populations start near
$\gamma=0.5$ and run \ScGens{} generations, with \ScPops{} populations per rule and 20 per
point in the sweeps. The values are chosen so that the gamble is worse on average for
everyone, and only the shape of the reward can favor it.

\section{Two illustrative models}\label{app:models}

\para{Concession.} A dominant tolerates $j$ subordinates, each of whom takes a share $\xi$ of
what it holds, and with them faces rivals of total strength $H$. It prevails with
probability $(1+j)^{\theta}/((1+j)^{\theta}+H^{\theta})$, where $\theta=1$ when fights are
duels and $\theta=2$ when attack can be massed, the linear and square laws of Lanchester
\cite{lanchester1916}. It chooses $j$ to maximize $(1-j\xi)$ times that probability.
Figure~\ref{fig:evo}a uses $\xi=0.05$. As $H$ grows, the share it keeps, $1-j\xi$, tends to
$(1+\xi)/(1+\theta)$. The model is in the tradition of concession models of reproductive
skew \cite{vehrencamp1983}, and it is an illustration with assumed parameters.

\para{Validity.} Figure~\ref{fig:evo}b evaluates the two-facet formula of
\S\ref{sec:evolution} for one assessor, with signal variance $v_s=0.5$, shared error
$v_0=0.1$, context error $v_c=2$ and occasion error $v_e=3$. Repeating one context $n$ times
gives $0.5/(2.6+3/n)$, with a ceiling of 19\%. One trial in each of $n$ contexts gives
$0.5/(0.6+5/n)$, with a ceiling of 83\%.

\section{What a deployment should publish}\label{app:record}

The conditions of \S\ref{sec:conditions} are claims about a particular system, and they can
be checked only if they are written down. A deployment that confers authority on an AI
system should keep a versioned record with at least the entries below. ``Unknown'' is an
admissible entry, and a more useful one than an invented small number.

\begin{enumerate}[itemsep=1pt,topsep=3pt]
\item \emph{Purpose and owners}: the task, the people affected, who grants authority, and who
reviews independently.
\item \emph{Scope}: kinds of action, destinations, data classes, credentials, shared
services, and the channels through which the system addresses people.
\item \emph{Identity}: model, runtime, tool and policy versions; memory and retention rules.
\item \emph{Authority}: limits per operation and in aggregate across agents and copies;
expiry; maximum concurrency and duration; how revocation works.
\item \emph{Properties}: forbidden effects, the invariants enforced, the trusted base, and
known omissions.
\item \emph{Usefulness}: the legitimate tasks that must still complete, and the rate of false
blocks.
\item \emph{The four numbers}: depth $k$ on the shallowest known route; the joint miss rate
under adaptive attack, with the attacker and the unit of trial; the search counts $n$ and
$n_{1}$ with the assumed advantage $\alpha$; and detection, decision and halt latency in the
tail.
\item \emph{Coverage in operation}: the share of activity each gate actually sees, as
distinct from the share it is meant to see.
\item \emph{Transfer}: which conditions of deployment match the evaluation, and which
opponents and circumstances are excluded.
\item \emph{Renewal}: which changes void which evidence, what reduced scope applies
meanwhile, and who decides.
\end{enumerate}
No single score replaces the record. A high task score does not repair an omitted channel,
an invariant does not certify the ethics of its specification, and a low failure rate per
episode does not justify unbounded volume.

Evidence has a validity domain and not only a date. Table~\ref{tab:renewal} lists, for the
commonest changes, what can stand and what must be established again before authority
continues.

\begin{table*}[t]
\centering
\footnotesize
\setlength{\tabcolsep}{5pt}
\renewcommand{\arraystretch}{1.16}
\begin{tabularx}{\textwidth}{@{}P{1.35in}LL@{}}
\toprule
\textbf{What changed} & \textbf{Evidence that can stand} & \textbf{What must be established again}\\
\midrule
The model or its weights & An invariant proved for every policy, under the same interface and dynamics & Behavioral and monitor evidence, including $J_{k}$ and $\alpha$; that tools and channels did not change with it\\
Copies, speed or duration & Invariants whose model of actors and scheduling already covers the change & Exposure: the number of attempts and opportunities, with one ledger preserved across all of them\\
Runtime, tools or credentials & Abstract properties and their proofs & That the implementation refines the model, and mediation for every new effect\\
Restart or renewed grant & Evidence whose premises persist through the transition & That consumption and identifier history were preserved; liveness and version checked again at commit\\
The specification, or the interest it protects & Evidence for properties the change does not touch & The new predicates, validated independently; every earlier check and useful task they affect\\
A new channel to people & Only claims that already covered that channel & What the channel can induce people to do; the term stays unknown where mediation is unproved\\
\bottomrule
\end{tabularx}
\caption{What survives a change. Keeping a result that is still valid and withholding an
extrapolation that is not are compatible. This is a proposed protocol and not an observed
practice.}
\label{tab:renewal}
\end{table*}

\normalsize
\begin{thebibliography}{199}

\bibitem{holloway2004} R. Holloway, D. Broadfield, M. Yuan. \emph{The Human Fossil Record, Vol. 3: Brain Endocasts}. Wiley, 2004.
\bibitem{du2018} A. Du et al. Pattern and process in hominin brain size evolution are scale-dependent. \emph{Proceedings of the Royal Society B} 285:20172738, 2018.
\bibitem{berger2015} L. Berger et al. \emph{Homo naledi}, a new species of the genus \emph{Homo} from the Dinaledi Chamber, South Africa. \emph{eLife} 4:e09560, 2015.
\bibitem{hawks2017} J. Hawks et al. New fossil remains of \emph{Homo naledi} from the Lesedi Chamber, South Africa. \emph{eLife} 6:e24232, 2017.
\bibitem{kubo2013} D. Kubo, R. Kono, Y. Kaifu. Brain size of \emph{Homo floresiensis} and its evolutionary implications. \emph{Proceedings of the Royal Society B} 280:20130338, 2013.
\bibitem{gingerich2022} P. Gingerich. Pattern and rate in the Plio-Pleistocene evolution of modern human brain size. \emph{Scientific Reports} 12:11216, 2022.
\bibitem{taylor2026pbo} J. Taylor. Performance before outcomes: early assessment, what women and men are chosen for, and the making of a general-purpose brain. Preprint, November 10, 2026.
\bibitem{schoen2025} B. Schoen et al. Stress testing deliberative alignment for anti-scheming training. arXiv:2509.15541, 2025.
\bibitem{needham2025} J. Needham, G. Edkins, G. Pimpale, H. Bartsch, M. Hobbhahn. Large language models often know when they are being evaluated. arXiv:2505.23836, 2025.
\bibitem{bengio2026iasr} Y. Bengio et al. International AI Safety Report 2026. February 2026.
\bibitem{openai2026astra} OpenAI. GPT-6 Astra system card and safety overview. September 2026.
\bibitem{anthropic2026judges} A. Lynch, J. Hughes, A. Serrano, R. Kirk, S. R. Bowman. Agentic misalignment in summer 2026. Anthropic Alignment Science, July 13, 2026.
\bibitem{boehm1993} C. Boehm. Egalitarian behavior and reverse dominance hierarchy. \emph{Current Anthropology} 34:227--254, 1993.
\bibitem{boehm1999} C. Boehm. \emph{Hierarchy in the Forest}. Harvard University Press, 1999.
\bibitem{woodburn1982} J. Woodburn. Egalitarian societies. \emph{Man} 17:431--451, 1982.
\bibitem{bostrom2012} N. Bostrom. The superintelligent will. \emph{Minds and Machines} 22:71--85, 2012.
\bibitem{omohundro2008} S. Omohundro. The basic AI drives. \emph{Proc. AGI}, 2008.
\bibitem{turner2021} A. Turner, L. Smith, R. Shah, A. Critch, P. Tadepalli. Optimal policies tend to seek power. \emph{NeurIPS}, 2021.
\bibitem{amodei2016} D. Amodei, C. Olah, J. Steinhardt, P. Christiano, J. Schulman, D. Man\'e. Concrete problems in AI safety. arXiv:1606.06565, 2016.
\bibitem{hubinger2019} E. Hubinger, C. van Merwijk, V. Mikulik, J. Skalse, S. Garrabrant. Risks from learned optimization in advanced machine learning systems. arXiv:1906.01820, 2019.
\bibitem{langosco2022} L. Langosco, J. Koch, L. Sharkey, J. Pfau, D. Krueger. Goal misgeneralization in deep reinforcement learning. \emph{ICML}, 2022.
\bibitem{soares2015} N. Soares, B. Fallenstein, E. Yudkowsky, S. Armstrong. Corrigibility. \emph{AAAI Workshop on AI and Ethics}, 2015.
\bibitem{hadfield2017} D. Hadfield-Menell, A. Dragan, P. Abbeel, S. Russell. The off-switch game. \emph{IJCAI}, 2017.
\bibitem{hudson2026} R. Hudson. Corrigibility transformation: constructing goals that accept updates. \emph{ICML}, PMLR 306, 2026.
\bibitem{carlsmith2022} J. Carlsmith. Is power-seeking AI an existential risk? arXiv:2206.13353, 2022.
\bibitem{hubinger2024} E. Hubinger et al. Sleeper agents: training deceptive LLMs that persist through safety training. arXiv:2401.05566, 2024.
\bibitem{greenblatt2024faking} R. Greenblatt et al. Alignment faking in large language models. arXiv:2412.14093, 2024.
\bibitem{meinke2024} A. Meinke et al. Frontier models are capable of in-context scheming. arXiv:2412.04984, 2024.
\bibitem{vanderweij2025} T. van der Weij, F. Hofst\"atter, O. Jaffe, S. Brown, F. Ward. AI sandbagging: language models can strategically underperform on evaluations. \emph{ICLR}, 2025.
\bibitem{sheshadri2025} A. Sheshadri et al. Why do some language models fake alignment while others don't? arXiv:2506.18032, 2025.
\bibitem{summerfield2025} C. Summerfield et al. Lessons from a chimp: AI ``scheming'' and the quest for ape language. arXiv:2507.03409, 2025.
\bibitem{baek2026} D. Baek, X. Li, A. Gupta, T. Mahbub, K. Shi, M. Tegmark, S. Feng. Sycophancy towards researchers drives performative misalignment. arXiv:2606.08629, 2026.
\bibitem{metr2026} H. Wijk, A. Cotra, R. Greenblatt. Brief independent investigation of agents' behavior, reasoning and collaboration in the OpenAI / Hugging Face hacking incident. METR and Redwood Research, August 26, 2026.
\bibitem{anthropic2026cyber} Anthropic. An alignment assessment of recent cybersecurity incidents. September 9, 2026.
\bibitem{anthropic2026aisi} Anthropic. An alignment assessment of the UK AISI cybertesting incident. October 23, 2026.
\bibitem{aisi2026incident} UK AI Security Institute. Investigating alignment failures in AISI's incident. October 23, 2026.
\bibitem{metr2026audit} D. Rein, A. Kaufman. Red-teaming OpenAI's internal agent monitoring systems. METR, October 22, 2026 (updated October 29).
\bibitem{marks2025} S. Marks et al. Auditing language models for hidden objectives. arXiv:2503.10965, 2025.
\bibitem{greenblatt2024control} R. Greenblatt, B. Shlegeris, K. Sachan, F. Roger. AI control: improving safety despite intentional subversion. \emph{ICML}, 2024.
\bibitem{moon2026} R. Moon, L. R. Varshney. Containment verification: AI safety guarantees independent of alignment. arXiv:2605.09045, 2026.
\bibitem{chen2026aar} Y.-H. Chen, J. Wen, J. H. Kirchner. Automated researchers can mitigate well-characterized alignment failures. arXiv:2608.28945, 2026.
\bibitem{balesni2024} M. Balesni et al. Towards evaluations-based safety cases for AI scheming. arXiv:2411.03336, 2024.
\bibitem{dalrymple2024} D. Dalrymple et al. Towards guaranteed safe AI: a framework for ensuring robust and reliable AI systems. arXiv:2405.06624, 2024.
\bibitem{engels2025} J. Engels, D. Baek, S. Kantamneni, M. Tegmark. Scaling laws for scalable oversight. arXiv:2504.18530, 2025.
\bibitem{carro2026} M. V. Carro et al. Humans correct their judgments through debate, but weaker AI models do not. arXiv:2610.28539, 2026.
\bibitem{terekhov2025} M. Terekhov et al. Adaptive attacks on trusted monitors subvert AI control protocols. arXiv:2510.09462, 2025.
\bibitem{makins2026} O. Makins, O. Angelini, Z. Shams, M. Phuong. Multi-agent AI control: distributed attacks hamper per-instance monitors. arXiv:2607.07368, 2026.
\bibitem{bowkis2026} A. Bowkis, M. D. Buhl, J. Pfau, G. Irving. Automated alignment is harder than you think. arXiv:2605.06390, 2026.
\bibitem{baker2025} B. Baker et al. Monitoring reasoning models for misbehavior and the risks of promoting obfuscation. arXiv:2503.11926, 2025.
\bibitem{korbak2025} T. Korbak et al. Chain of thought monitorability: a new and fragile opportunity for AI safety. arXiv:2507.11473, 2025.
\bibitem{khan2024} A. Khan et al. Debating with more persuasive LLMs leads to more truthful answers. \emph{ICML}, 2024.
\bibitem{irving2018} G. Irving, P. Christiano, D. Amodei. AI safety via debate. arXiv:1805.00899, 2018.
\bibitem{burns2024} C. Burns et al. Weak-to-strong generalization: eliciting strong capabilities with weak supervision. \emph{ICML}, 2024.
\bibitem{klein2009} G. Klein et al. seL4: formal verification of an OS kernel. \emph{SOSP}, 2009.
\bibitem{bengio2025scientist} Y. Bengio et al. Superintelligent agents pose catastrophic risks: can Scientist AI offer a safer path? arXiv:2502.15657, 2025.
\bibitem{vermeer2025} M. J. D. Vermeer, E. Lathrop, A. Moon. On the extinction risk from artificial intelligence. RAND RR-A3034-1, 2025.
\bibitem{narayanan2025} A. Narayanan, S. Kapoor. AI as normal technology. Knight First Amendment Institute, 2025.
\bibitem{hagele2026} A. H\"agele et al. The hot mess of AI. arXiv:2601.23045, 2026.
\bibitem{hong2026} S. Z. Hong et al. Measuring mid-2025 LLM-assistance on novice performance in biology. arXiv:2602.16703, 2026.
\bibitem{grace2025} K. Grace et al. Thousands of AI authors on the future of AI. \emph{JAIR} 84, 2025.
\bibitem{kulveit2025} J. Kulveit et al. Position: humanity faces existential risk from gradual disempowerment. \emph{ICML}, 2025.
\bibitem{yudkowsky2025} E. Yudkowsky, N. Soares. \emph{If Anyone Builds It, Everyone Dies}. Little, Brown, 2025.
\bibitem{jones2025} E. Jones et al. Forecasting rare language model behaviors. arXiv:2502.16797, 2025.
\bibitem{wu2024} G. Wu, J. Hilton. Estimating the probabilities of rare outputs in language models. arXiv:2410.13211, 2024.
\bibitem{anderson1972} J. P. Anderson. \emph{Computer Security Technology Planning Study}. ESD-TR-73-51, US Air Force, 1972.
\bibitem{saltzer1975} J. Saltzer, M. Schroeder. The protection of information in computer systems. \emph{Proceedings of the IEEE} 63:1278--1308, 1975.
\bibitem{schneider2000} F. Schneider. Enforceable security policies. \emph{ACM Transactions on Information and System Security} 3:30--50, 2000.
\bibitem{alshiekh2018} M. Alshiekh et al. Safe reinforcement learning via shielding. \emph{AAAI}, 2018.
\bibitem{fleming1975} K. N. Fleming. A reliability model for common mode failures in redundant safety systems. Report GA-A13284, General Atomic, 1975.
\bibitem{mosleh1988} A. Mosleh et al. \emph{Procedures for Treating Common Cause Failures in Safety and Reliability Studies}. NUREG/CR-4780, 1988.
\bibitem{leveson2011} N. Leveson. \emph{Engineering a Safer World}. MIT Press, 2011.
\bibitem{ord2020} T. Ord. \emph{The Precipice}. Bloomsbury, 2020.
\bibitem{kwa2025} T. Kwa et al. Measuring AI ability to complete long tasks. arXiv:2503.14499, 2025.
\bibitem{urbina2022} F. Urbina, F. Lentzos, C. Invernizzi, S. Ekins. Dual use of artificial-intelligence-powered drug discovery. \emph{Nature Machine Intelligence} 4:189--191, 2022.
\bibitem{schneier1999} B. Schneier. Attack trees. \emph{Dr. Dobb's Journal} 24(12):21--29, 1999.
\bibitem{sheyner2002} O. Sheyner, J. Haines, S. Jha, R. Lippmann, J. Wing. Automated generation and analysis of attack graphs. \emph{IEEE Symposium on Security and Privacy}, 2002.
\bibitem{wash1400} US Nuclear Regulatory Commission. \emph{Reactor Safety Study}. WASH-1400 (NUREG-75/014), 1975.
\bibitem{menger1927} K. Menger. Zur allgemeinen Kurventheorie. \emph{Fundamenta Mathematicae} 10:96--115, 1927.
\bibitem{ford1956} L. Ford, D. Fulkerson. Maximal flow through a network. \emph{Canadian Journal of Mathematics} 8:399--404, 1956.
\bibitem{robacker1956} J. Robacker. Min-max theorems on shortest chains and disjoint cuts of a network. RAND RM-1660, 1956.
\bibitem{hhs2023} US Department of Health and Human Services. \emph{Screening Framework Guidance for Providers and Users of Synthetic Nucleic Acids}. 2023.
\bibitem{esvelt2022} K. Esvelt. Delay, detect, defend: preparing for a future in which thousands can release new pandemics. Geneva Centre for Security Policy, 2022.
\bibitem{berg1975} P. Berg, D. Baltimore, S. Brenner, R. Roblin, M. Singer. Summary statement of the Asilomar conference on recombinant DNA molecules. \emph{PNAS} 72:1981--1984, 1975.
\bibitem{adamala2024} K. Adamala et al. Confronting risks of mirror life. \emph{Science} 386:1351--1353, 2024.
\bibitem{popper1959} K. Popper. \emph{The Logic of Scientific Discovery}. Hutchinson, 1959.
\bibitem{good1953} I. J. Good. The population frequencies of species and the estimation of population parameters. \emph{Biometrika} 40:237--264, 1953.
\bibitem{bohme2018} M. B\"ohme. STADS: software testing as species discovery. \emph{ACM Transactions on Software Engineering and Methodology} 27(2), 2018.
\bibitem{mcallester2000} D. McAllester, R. Schapire. On the convergence rate of Good--Turing estimators. \emph{COLT}, 2000.
\bibitem{hanley1983} J. Hanley, A. Lippman-Hand. If nothing goes wrong, is everything all right? \emph{JAMA} 249:1743--1745, 1983.
\bibitem{clopper1934} C. Clopper, E. Pearson. The use of confidence or fiducial limits illustrated in the case of the binomial. \emph{Biometrika} 26:404--413, 1934.
\bibitem{greenblatt2024password} R. Greenblatt, F. Roger, D. Krasheninnikov, D. Krueger. Stress-testing capability elicitation with password-locked models. arXiv:2405.19550, 2024.
\bibitem{bostrom2011} N. Bostrom. Information hazards: a typology of potential harms from knowledge. \emph{Review of Contemporary Philosophy} 10:44--79, 2011.
\bibitem{bostrom2019} N. Bostrom. The vulnerable world hypothesis. \emph{Global Policy} 10:455--476, 2019.
\bibitem{armstrong2012} S. Armstrong, A. Sandberg, N. Bostrom. Thinking inside the box: controlling and using an Oracle AI. \emph{Minds and Machines} 22:299--324, 2012.
\bibitem{betley2025} J. Betley et al. Emergent misalignment: narrow finetuning can produce broadly misaligned LLMs. \emph{ICML}, 2025.
\bibitem{shamir1979} A. Shamir. How to share a secret. \emph{Communications of the ACM} 22:612--613, 1979.
\bibitem{lande1983} R. Lande, S. Arnold. The measurement of selection on correlated characters. \emph{Evolution} 37:1210--1226, 1983.
\bibitem{trut1999} L. Trut. Early canid domestication: the farm-fox experiment. \emph{American Scientist} 87:160--169, 1999.
\bibitem{price1970} G. Price. Selection and covariance. \emph{Nature} 227:520--521, 1970.
\bibitem{plavcan2012} J. M. Plavcan. Sexual size dimorphism, canine dimorphism, and male--male competition in primates: where do humans fit in? \emph{Human Nature} 23:45--67, 2012.
\bibitem{lazear1981} E. Lazear, S. Rosen. Rank-order tournaments as optimum labor contracts. \emph{Journal of Political Economy} 89:841--864, 1981.
\bibitem{hvide2002} H. Hvide. Tournament rewards and risk taking. \emph{Journal of Labor Economics} 20:877--898, 2002.
\bibitem{wilson1985} M. Wilson, M. Daly. Competitiveness, risk taking, and violence: the young male syndrome. \emph{Ethology and Sociobiology} 6:59--73, 1985.
\bibitem{bondarenko2025} A. Bondarenko, D. Volk, D. Volkov, J. Ladish. Demonstrating specification gaming in reasoning models. arXiv:2502.13295, 2025.
\bibitem{henrich2001} J. Henrich, F. Gil-White. The evolution of prestige. \emph{Evolution and Human Behavior} 22:165--196, 2001.
\bibitem{cheng2013} J. Cheng, J. Tracy, T. Foulsham, A. Kingstone, J. Henrich. Two ways to the top. \emph{Journal of Personality and Social Psychology} 104:103--125, 2013.
\bibitem{gavrilets2012} S. Gavrilets. On the evolutionary origins of the egalitarian syndrome. \emph{PNAS} 109:14069--14074, 2012.
\bibitem{wrangham2019} R. Wrangham. \emph{The Goodness Paradox}. Pantheon, 2019.
\bibitem{vehrencamp1983} S. Vehrencamp. A model for the evolution of despotic versus egalitarian societies. \emph{Animal Behaviour} 31:667--682, 1983.
\bibitem{bingham1999} P. Bingham. Human uniqueness: a general theory. \emph{Quarterly Review of Biology} 74:133--169, 1999.
\bibitem{borgerhoff2009} M. Borgerhoff Mulder et al. Intergenerational wealth transmission and the dynamics of inequality in small-scale societies. \emph{Science} 326:682--688, 2009.
\bibitem{betzig1986} L. Betzig. \emph{Despotism and Differential Reproduction}. Aldine, 1986.
\bibitem{carneiro1970} R. Carneiro. A theory of the origin of the state. \emph{Science} 169:733--738, 1970.
\bibitem{smuts1993} B. Smuts, R. Smuts. Male aggression and sexual coercion of females in nonhuman primates and other mammals. \emph{Advances in the Study of Behavior} 22:1--63, 1993.
\bibitem{hrdy1979} S. Hrdy. Infanticide among animals: a review, classification, and examination of the implications for the reproductive strategies of females. \emph{Ethology and Sociobiology} 1:13--40, 1979.
\bibitem{chen2025persona} R. Chen, A. Arditi, H. Sleight, O. Evans, J. Lindsey. Persona vectors: monitoring and controlling character traits in language models. arXiv:2507.21509, 2025.
\bibitem{hawkes1998} K. Hawkes, J. O'Connell, N. Blurton Jones, H. Alvarez, E. Charnov. Grandmothering, menopause, and the evolution of human life histories. \emph{PNAS} 95:1336--1339, 1998.
\bibitem{lahdenpera2004} M. Lahdenper\"a, V. Lummaa, S. Helle, M. Tremblay, A. Russell. Fitness benefits of prolonged post-reproductive lifespan in women. \emph{Nature} 428:178--181, 2004.
\bibitem{hrdy2009} S. Hrdy. \emph{Mothers and Others}. Harvard University Press, 2009.
\bibitem{chapman2019} S. Chapman, J. Pettay, V. Lummaa, M. Lahdenper\"a. Limits to fitness benefits of prolonged post-reproductive lifespan in women. \emph{Current Biology} 29:645--650, 2019.
\bibitem{glowacki2015} L. Glowacki, R. Wrangham. Warfare and reproductive success in a tribal population. \emph{PNAS} 112:348--353, 2015.
\bibitem{beckerman2009} S. Beckerman et al. Life histories, blood revenge, and reproductive success among the Waorani of Ecuador. \emph{PNAS} 106:8134--8139, 2009.
\bibitem{walker2011} R. Walker, K. Hill, M. Flinn, R. Ellsworth. Evolutionary history of hunter-gatherer marriage practices. \emph{PLoS ONE} 6:e19066, 2011.
\bibitem{apostolou2007} M. Apostolou. Sexual selection under parental choice: the role of parents in the evolution of human mating. \emph{Evolution and Human Behavior} 28:403--409, 2007.
\bibitem{cronbach1972} L. Cronbach, G. Gleser, H. Nanda, N. Rajaratnam. \emph{The Dependability of Behavioral Measurements: Theory of Generalizability for Scores and Profiles}. Wiley, 1972.
\bibitem{marlowe2004} F. Marlowe. Marital residence among foragers. \emph{Current Anthropology} 45:277--284, 2004.
\bibitem{papernot2016} N. Papernot, P. McDaniel, I. Goodfellow. Transferability in machine learning: from phenomena to black-box attacks using adversarial samples. arXiv:1605.07277, 2016.
\bibitem{tramer2017} F. Tram\`er, N. Papernot, I. Goodfellow, D. Boneh, P. McDaniel. The space of transferable adversarial examples. arXiv:1704.03453, 2017.
\bibitem{lampson1973} B. Lampson. A note on the confinement problem. \emph{Communications of the ACM} 16:613--615, 1973.
\bibitem{kogui2026} K\'aggaba (Kogui) people; Mamo Vangelio Moscote Alberto and Saka Rosa Mamatac\'an Sauna, Ezuama Mayor Kuamaka. An urgent message on artificial intelligence from the Kogui people. Sierra Nevada de Gonawindua, Colombia. Undated; posted September 2026. \url{https://movilizatorio.org/wp-content/uploads/2026/09/Gonawindua-Kogui-AI-Statement-.pdf}
\bibitem{herlihy1990} M. Herlihy, J. Wing. Linearizability: a correctness condition for concurrent objects. \emph{ACM Transactions on Programming Languages and Systems} 12:463--492, 1990.
\bibitem{lanchester1916} F. Lanchester. \emph{Aircraft in Warfare: The Dawn of the Fourth Arm}. Constable, 1916.

\end{thebibliography}
\end{document}
